% ============================================================
% Master's thesis skeleton — English version, compliant with
% WIiT PK guidelines (Załącznik do Uchwały Kolegium WIiT
% z dn. 25.09.2023 r. nr 74/2023)
% ============================================================
% Formatting per the guidelines (language-independent):
%  - A4, portrait, double-sided printing (twoside)
%  - font: Times New Roman 12pt (newtxtext as the equivalent —
%    mathptmx was tested and found to silently drop several
%    Polish diacritics under T1 encoding; newtxtext does not)
%  - line spacing: 1.5
%  - margins: top 2.5 cm / bottom 2.5 cm / inner 3.5 cm / outer 2.5 cm
%  - paragraph indent: 1.25 cm
%  - page numbering: from the table of contents (p. 2), 10pt, centered
%  - title page: page number must not be visible
%  - NOTE: per the guidelines, an English-language thesis must still
%    include a Polish-language "Streszczenie" (abstract) — kept below,
%    typeset with \selectlanguage{polish} so Polish hyphenation/diacritics
%    are handled correctly even though the rest of the document is English.
% ============================================================

\documentclass[12pt,a4paper,twoside]{report}

% --- Encoding and language (English as main, Polish available for the
%     mandatory Polish abstract) ---
\usepackage[utf8]{inputenc}
\usepackage[T1]{fontenc}
\usepackage[polish,english]{babel}

% --- Font (Times New Roman equivalent) ---
% newtxtext confirmed by test compilation to render Polish diacritics
% correctly under T1 + babel — needed here for the mandatory Polish
% "Streszczenie" section even though the rest of the thesis is in English.
\usepackage{newtxtext}

% --- Margins per the guidelines (mirrored for double-sided printing) ---
\usepackage[a4paper,twoside,
  top=2.5cm, bottom=2.5cm,
  inner=3.5cm, outer=2.5cm]{geometry}

% --- 1.5 line spacing ---
\usepackage{setspace}
\onehalfspacing

% --- Let pages end at their natural height instead of stretching
%     vertical glue to fill every page exactly (report class default
%     is \flushbottom) --- the stretching was exaggerating the gaps
%     around tables/figures on shorter pages, making them look
%     stranded in the middle of mostly-blank pages. ---
\raggedbottom

% --- Paragraph indent ---
\setlength{\parindent}{1.25cm}

% --- Header/footer: centered page number, 10pt ---
\usepackage{fancyhdr}
\pagestyle{fancy}
\fancyhf{}
\renewcommand{\headrulewidth}{0pt}
\fancyfoot[C]{\fontsize{10pt}{12pt}\selectfont\thepage}

% --- Figures, tables, references ---
\usepackage{graphicx}
\usepackage{booktabs}
% Wrapping p{} columns for Table 4.1, whose long cell text (e.g.
% "entropy, chi-square, avalanche, correlation") overflowed the right
% margin under plain non-wrapping l columns.
\usepackage{array}
\usepackage{amsmath}
\usepackage{amssymb}   % \mathbb, used in Section 3.3
\usepackage{hyperref}
% hypertexnames=false suppresses a pre-existing duplicate-anchor warning that
% hyperref emits for every float in this document; harmless, but noisy.
\hypersetup{colorlinks=true, linkcolor=black, citecolor=black, urlcolor=blue,
  hypertexnames=false}

\usepackage{tikz}
\usetikzlibrary{arrows.meta, positioning, calc, decorations.pathreplacing}

% --- Forces the AES-GCM figure below to stay in Section 3.5, instead of
%     floating past the Section 3.6 heading when it does not fit on the
%     current page (observed with plain [htbp] placement) ---
\usepackage{float}

% --- Chapter 7 data visualizations (histogram, heatmap, bar chart) ---
\usepackage{xcolor}
\usepackage{pgfplots}
\pgfplotsset{compat=1.18}
% Sequential blue, matches this project's validated categorical palette
% (used elsewhere for M6's Artifact figures) reused here for consistency.
\definecolor{seqblue}{HTML}{2a78d6}

% --- Source-code listings (implementation snippets, Chapters 3, 5, 6) ---
\usepackage{listings}
\lstdefinestyle{pythonstyle}{
  language=Python,
  basicstyle=\ttfamily\small,
  keywordstyle=\bfseries,
  commentstyle=\itshape\color{gray!60!black},
  stringstyle=\color{seqblue!70!black},
  numbers=left,
  numberstyle=\tiny\color{gray},
  numbersep=8pt,
  frame=single,
  rulecolor=\color{gray!40},
  breaklines=true,
  showstringspaces=false,
  tabsize=4,
  xleftmargin=1.25cm,
  framexleftmargin=5pt,
}
\lstset{style=pythonstyle}
% Listings are numbered and captioned throughout, so the document also
% carries a List of Listings alongside the figure and table lists.
\renewcommand{\lstlistlistingname}{List of Listings}

% --- Equation numbering per chapter ---
\numberwithin{equation}{chapter}

\begin{document}

% ============================================================
% TITLE PAGE
% ------------------------------------------------------------
% Kept in its own file (title_page.tex). Per the WIiT PK
% guidelines the title page must carry no visible page number,
% so that file is responsible for \thispagestyle{empty} (or a
% titlepage environment, which applies it automatically).
% ============================================================
\input{title_page}

% ============================================================
% STRESZCZENIE (PL, mandatory even in an English thesis) / ABSTRACT (EN)
% ============================================================
\selectlanguage{polish}
\chapter*{Streszczenie}
\addcontentsline{toc}{chapter}{Streszczenie}

Niniejsza praca bada, czy transformacja kodowania DNA pozostawia
mierzalny odcisk statystyczny na danych o niskiej entropii (tekst, JSON)
w odróżnieniu od danych o wysokiej entropii (dane losowe), oraz czy
efekt ten zanika po nałożeniu klasycznego szyfrowania (AES-GCM). Zbadano
dwa potoki przetwarzania (z transformacją DNA i bez niej, jako kontrolę)
na zamrożonym korpusie testowym, mierząc entropię Shannona, statystykę
chi-kwadrat, współczynnik kompresji, korelację, efekt lawinowy oraz
narzut rozmiaru danych. Wyniki wskazują na niewielki sygnał statystyczny
na etapie samej transformacji DNA, który zanika całkowicie po
zastosowaniu szyfrowania klasycznego. Wnioskiem pracy jest to, że
transformacja w domenie DNA tego rodzaju nie stanowi zamiennika
szyfrowania klasycznego: jej własny wkład statystyczny jest mierzalny
przed szyfrowaniem, lecz nie przeżywa nałożenia AES-GCM.

\textbf{Słowa kluczowe:} kodowanie DNA, kryptografia lekka, IoT,
entropia Shannona, AES-GCM, analiza statystyczna

\selectlanguage{english}
\chapter*{Abstract}
\addcontentsline{toc}{chapter}{Abstract}

This thesis investigates whether a DNA-encoding transform leaves a
measurable statistical fingerprint on low-entropy data (text, JSON) as
opposed to high-entropy data (random), and whether that effect
disappears once classical encryption (AES-GCM) is layered on top. Two
processing pipelines (with and without the DNA transform, as a control)
were evaluated against a frozen test corpus, measuring Shannon entropy,
chi-square statistic, compression ratio, correlation, avalanche effect,
and size overhead. Results indicate a small statistical
signal at the DNA-transform stage alone, which disappears entirely once
classical encryption is applied.

\textbf{Keywords:} DNA encoding, lightweight cryptography, IoT, Shannon
entropy, AES-GCM, statistical analysis

% ============================================================
% TABLE OF CONTENTS (page 2, page number visible)
% ============================================================
\tableofcontents

% ============================================================
% LIST OF ABBREVIATIONS
% ============================================================
\chapter*{List of Abbreviations}
\addcontentsline{toc}{chapter}{List of Abbreviations}

\begin{description}
  \item[AEAD] Authenticated Encryption with Associated Data
  \item[AES-GCM] Advanced Encryption Standard, Galois/Counter Mode
  \item[ECC] Elliptic Curve Cryptography
  \item[DNA-LWCS] DNA-based Lightweight Cryptography System
  \item[IoT] Internet of Things
  \item[IV/nonce] Initialization Vector: a single-use value used in encryption
  \item[NIST] National Institute of Standards and Technology
  \item[UACI/NPCR] standard metrics used in image-encryption analysis
\end{description}

% ============================================================
% CHAPTER 1 --- INTRODUCTION
% ============================================================
\chapter{Introduction}
\label{ch:intro}

% No figures/tables here --- narrative only: motivation, problem
% justification, literature context.

By the end of 2025, industry estimates put the number of active IoT
devices worldwide at roughly 21~billion, up about 14\% over the
previous year~\cite{iotanalytics2025}. IoT-targeted attacks surged by
more than 100\% year-over-year in the first half of
2024~\cite{sonicwall2024}, and one industry report puts the sustained
volume at over 820{,}000 attack attempts recorded
daily~\cite{stationx2026}. Much of this traffic remains unencrypted: a
survey of 1.2~million IoT devices across enterprise and healthcare
networks found 98\% of their traffic was unencrypted~\cite{unit42-2020},
a figure still widely recycled today, years after that original survey.

Most IoT devices operate under tight resource constraints: limited CPU,
small memory, and a restricted energy budget~\cite{zinabu2025}. Security
therefore has to fit inside that budget, or it is left out entirely.
Conventional cryptographic standards
were not designed with such constraints in mind: AES, RSA, and ECC alike
are frequently too resource-intensive for the smallest devices, with AES
in particular exhibiting lower power efficiency and a larger memory
footprint than ciphers purpose-built for constrained
hardware~\cite{zinabu2025}. This gap motivated a dedicated field of
lightweight cryptography, aimed at preserving security while minimising
computation, memory use, and energy
consumption~\cite{thakor2021}. That effort was consolidated in 2023, when
the U.S. National Institute of Standards and Technology (NIST) selected
the Ascon family as the basis for its lightweight cryptography standard
for constrained devices, later published as NIST~SP~800-232~\cite{nist800232}.

DNA encoding is a research direction that tries to combine both
strengths. It maps binary data onto sequences of nucleotides (A, T, C,
G), then applies a transform to this DNA-domain representation, either
before classical encryption or alongside it. Two papers form the
starting point for this thesis's design. Aqeel et al.~\cite{aqeel2025}
propose DNA-LWCS. This scheme derives key material from public DNA
sequences. It combines a DNA-domain substitution-and-transposition step
with Elliptic Curve Cryptography (ECC). The authors report good results
for encryption time, throughput, and energy use on image data.
Barman and Saha~\cite{dnaecc} take a similar approach: they encode data
as DNA sequences before applying ECC-based encryption. The motivation
given is the large number of possible DNA sequences and their apparent
randomness.

Both works share the same general idea. Classical encryption still
provides the main confidentiality guarantee. The DNA-domain step is
added as an extra, lightweight layer on top. Chapter~\ref{ch:foundations}
describes the DNA-domain transform used in this thesis and how it
differs from these two proposals. Chapter~\ref{ch:related} gives a full
architectural comparison.

A gap remains in the papers reviewed for this thesis: DNA-based cryptographic schemes are tested mainly on image data.
DNA-LWCS itself reports its results on three sample
photographs~\cite{aqeel2025}. It reports system-level metrics such as
encryption time, correlation, and energy use. Neither reviewed paper
tests whether the DNA-encoding step leaves a statistical fingerprint
that depends on the structure of the input data. Neither tests whether
such a fingerprint would survive once classical encryption is applied
on top, even though a large part of IoT traffic is structured text
and JSON: sensor readings, configuration messages, log records, and
binary payloads. This thesis targets that gap. It tests whether a
DNA-encoding transform's statistical behaviour differs across three data
classes: low-entropy text, low-entropy JSON, and high-entropy random or
binary data. It also compares an AES-GCM cipher with and without the
DNA-encoding step. This before/after comparison does not appear in the
two reviewed papers, where the DNA step is built into the encryption
pipeline and not tested as a separate stage.

One caveat regarding AES-GCM is that this thesis does not test whether DNA
encoding is "lightweight" in a performance sense. AES-GCM is used as a
standard, controlled baseline. Its role is to isolate the DNA
transform's own effect, not to claim that the combined scheme is
lightweight. Chapter~\ref{ch:limitations} explains why performance is
measured but deliberately not reported as a result of this thesis.

% ============================================================
% CHAPTER 2 --- AIM AND SCOPE OF THE THESIS
% ============================================================
\chapter{Aim and Scope of the Thesis}
\label{ch:aim}

% Concise --- no technical detail, no tables/figures here.

The aim of this thesis, and its research question, is to find out
whether a DNA-encoding transform produces a measurably different
statistical signature on low-entropy data (text, JSON) than on
high-entropy data (random or compressed), and whether that effect
disappears once classical encryption (AES-GCM) is layered on top.

This aim is stated as two testable hypotheses, each evaluated in
Chapter~\ref{ch:results} and answered in Chapter~\ref{ch:discussion}:

\begin{description}
\item[Hypothesis 1 (H1)] At the post-DNA-transform stage, the byte-value distribution
of low-entropy input (text, JSON) departs from uniformity by
measurably more than that of high-entropy input (random). H1 is
rejected if the chi-square statistics of the three data classes are
indistinguishable at that stage, given the between-sample spread
reported in Section~\ref{sec:results-repeated-sample}.

\item[Hypothesis 2 (H2)] After AES-GCM encryption, no such difference between data
classes remains, and the measurements do not depend on whether the DNA
step ran. H2 is rejected if any metric separates the data classes, or
separates Pipeline A from Pipeline B, at the post-classical stage.
\end{description}

Scope of the thesis covers three data classes (text, JSON,
random), two processing pipelines (DNA transform combined with AES-GCM,
and AES-GCM alone as a control), and three measurement stages (input
data, post-DNA-transform, post-classical-encryption).

The statistical testing used in this thesis does not replace formal
cryptanalysis. This caveat is adapted from NIST SP 800-22~\cite{nist80022}.
The scope leaves out the following on purpose: the full NIST STS test
suite, ECC/ECDH, image data, and energy-consumption profiling.
Chapter~\ref{ch:limitations} explains why.

The thesis is organized as follows: Chapter~\ref{ch:foundations}
presents the theoretical foundations of the mechanisms used;
Chapter~\ref{ch:related} discusses related work;
Chapter~\ref{ch:methodology} describes the research methodology;
Chapter~\ref{ch:implementation} covers the implementation;
Chapter~\ref{ch:results} presents the results obtained;
Chapter~\ref{ch:discussion} discusses the results;
Chapter~\ref{ch:limitations} covers the thesis's limitations; and
Chapter~\ref{ch:summary} summarises the contributions and outlines
future work.

% ============================================================
% CHAPTER 3 --- THEORETICAL FOUNDATIONS
% ============================================================
\chapter{Theoretical Foundations}
\label{ch:foundations}

\section{Bytes to DNA Sequences Transformation}
\label{sec:bytes-to-dna}
 Before any DNA-domain operation can be applied, the plaintext bytes must first be represented as nucleotides. Every byte consists of 8 bits; every 2 bits map to one nucleotide (A=00, T=01, C=10, G=11). A binary file becomes a sequence of A/T/C/G letters. This mapping is reversible, but does not bring any cryptographic properties on its own.

\begin{figure}[htbp]
\centering
\begin{tikzpicture}[
  bitbox/.style={draw, minimum width=0.75cm, minimum height=0.9cm, font=\ttfamily\large},
  nucbox/.style={draw, minimum width=1.5cm, minimum height=0.9cm, font=\ttfamily\bfseries\large, fill=gray!12},
  >=Stealth
]
  \node[bitbox] (b0) at (0.00,0) {1};
  \node[bitbox] (b1) at (0.75,0) {0};
  \node[bitbox] (b2) at (1.50,0) {1};
  \node[bitbox] (b3) at (2.25,0) {1};
  \node[bitbox] (b4) at (3.00,0) {0};
  \node[bitbox] (b5) at (3.75,0) {0};
  \node[bitbox] (b6) at (4.50,0) {0};
  \node[bitbox] (b7) at (5.25,0) {1};
  \node[anchor=east] at ($(b0.west)+(-0.3,0)$) {Byte:};

  \draw[decorate, decoration={brace, amplitude=4pt, mirror}]
    (b0.south west) -- (b1.south east) node[midway, yshift=-0.4cm, font=\scriptsize] {10};
  \draw[decorate, decoration={brace, amplitude=4pt, mirror}]
    (b2.south west) -- (b3.south east) node[midway, yshift=-0.4cm, font=\scriptsize] {11};
  \draw[decorate, decoration={brace, amplitude=4pt, mirror}]
    (b4.south west) -- (b5.south east) node[midway, yshift=-0.4cm, font=\scriptsize] {00};
  \draw[decorate, decoration={brace, amplitude=4pt, mirror}]
    (b6.south west) -- (b7.south east) node[midway, yshift=-0.4cm, font=\scriptsize] {01};

  \node[nucbox] (n0) at (0.375,-2.3) {C};
  \node[nucbox] (n1) at (1.875,-2.3) {G};
  \node[nucbox] (n2) at (3.375,-2.3) {A};
  \node[nucbox] (n3) at (4.875,-2.3) {T};
  \node[anchor=east] at ($(n0.west)+(-0.6,0)$) {DNA:};

  \draw[->] ($(b0)!0.5!(b1)$) -- (n0.north);
  \draw[->] ($(b2)!0.5!(b3)$) -- (n1.north);
  \draw[->] ($(b4)!0.5!(b5)$) -- (n2.north);
  \draw[->] ($(b6)!0.5!(b7)$) -- (n3.north);

  \node[anchor=west, font=\scriptsize] at (6.4,-2.3) {%
    \begin{tabular}{ll}
    A = 00 & C = 10 \\
    T = 01 & G = 11 \\
    \end{tabular}};
\end{tikzpicture}
\caption{Byte-to-nucleotide mapping: every 2 bits of the input byte map
to one of four nucleotides, using the fixed A/T/C/G codebook shown on
the right.}
\label{fig:byte-to-dna}
\end{figure}

\section{Why a Naive DNA Transform Fails}
The first idea was to derive a mask from a DNA sequence, for example by
hashing it, and combine it with the data using XOR. This is similar to
the keystream-and-XOR approach that AES-GCM uses. This idea was tried
during the implementation of this thesis, and it was rejected. A
cryptographically strong mask pushes every input toward the same,
near-maximal entropy, regardless of what the original data looked like.
This follows from a basic property of XOR: if $M$ is uniformly
distributed and independent of the plaintext $P$, then
$C = P \oplus M$ is uniformly distributed regardless of $P$'s own
distribution. As a result, the mask erases the exact
signal this thesis needs to measure. See Chapter~\ref{ch:implementation}
for the empirical details.

\section{The Chained Modulo-4 Design}
\label{sec:mod4chain}
The final transform design uses addition in the nucleotide domain, with
feedback:

\begin{equation}
c_i = \left( p_i + r_{i \bmod L} + c_{i-1} \right) \bmod 4
\label{eq:mod4}
\end{equation}

where $p_i$ is the value of the $i$-th nucleotide of the input data,
$r$ is the reference sequence of length $L$, and $c_{i-1}$ is the
previously generated ciphertext nucleotide ($c_{-1} = 0$). Each output
nucleotide depends on the entire history of the stream so far. This
removes the periodicity found in a simpler design, one that just tiles
the reference sequence without feedback.

Equation~\ref{eq:mod4} is implemented directly in
\texttt{src/dna\_transform.py}, as a single pass over the input holding
one accumulator that carries $c_{i-1}$ from each iteration to the next:
the feedback path drawn in Fig.~\ref{fig:mod4-chain}.

\begin{figure}[htbp]
\centering
\begin{tikzpicture}[
  node distance=1.0cm and 1.5cm,
  block/.style={draw, minimum width=2.1cm, minimum height=1cm, align=center, font=\small},
  sumnode/.style={draw, circle, minimum size=0.9cm, font=\Large},
  outnode/.style={draw, minimum width=2.1cm, minimum height=1cm, align=center, font=\small, fill=gray!12},
  >=Stealth
]
  \node[block] (p) {$p_i$\\[2pt]input nucleotide};
  \node[sumnode, right=of p] (sum) {$+$};
  \node[block, right=of sum] (mod) {$\bmod\ 4$};
  \node[outnode, right=of mod] (out) {$c_i$\\[2pt]output nucleotide};
  \node[block, below=1.6cm of sum] (prev) {$c_{i-1}$\\[2pt]previous output};

  \draw[->] (p) -- (sum);
  \draw[->] (mod) -- (out);
  \draw[->] (sum) -- (mod);
  \draw[->] (prev.north) -- (sum.south);

  \node[block, above=1.0cm of sum] (ref) {$r_{i \bmod L}$\\[2pt]reference sequence};
  \draw[->] (ref.south) -- (sum.north);

  \draw[->, dashed] (out.east) -- ++(1.2,0) |- node[near start, right, font=\scriptsize, align=left] {feeds forward as\\$c_{i-1}$ of step $i+1$} ++(0,-3.6) -| (prev.east);
\end{tikzpicture}
\caption{One step of the chained modulo-4 transform. Each output
nucleotide $c_i$ combines the input nucleotide $p_i$, the matching
reference-sequence nucleotide $r_{i \bmod L}$, and the previous step's
own output $c_{i-1}$. That output is then carried forward into the next
step. }
\label{fig:mod4-chain}
\end{figure}

\subsection*{Closed form: the transform is a prefix sum}
Equation~\ref{eq:mod4} is defined recursively, but the recursion
unrolls into a closed form. Substituting $c_{i-1}$ repeatedly, and
using $c_{-1} = 0$, gives

\begin{equation}
c_i = \left( \sum_{j=0}^{i} \left( p_j + r_{j \bmod L} \right) \right) \bmod 4 .
\label{eq:prefixsum}
\end{equation}

Each output nucleotide is therefore a running total, modulo 4, of every
input nucleotide and every reference nucleotide seen so far. Two
properties follow directly, and both are used later in this thesis.

First, the transform propagates a change in one direction only. If a
single input nucleotide $p_k$ changes by $\delta$, then by
Equation~\ref{eq:prefixsum} every output nucleotide $c_i$ with $i \ge k$
changes by exactly the same constant $\delta \bmod 4$, and every output
nucleotide with $i < k$ is unchanged. Section~\ref{sec:results-avalanche}
uses this to explain the measured avalanche figure.

Second, the reference sequence enters only through its own running
total. Writing $R_i = \sum_{j \le i} r_{j \bmod L} \bmod 4$ for that
accumulated mask, and $S = \sum_{j=0}^{L-1} r_j \bmod 4$ for the sum
over one full period of the reference sequence, we have
$R_{i+L} = (R_i + S) \bmod 4$. The accumulated mask therefore repeats
with period $L$ only when $S \equiv 0 \pmod 4$; otherwise its period is
$2L$ (when $S \equiv 2$) or $4L$ (when $S \equiv 1$ or $3$).
Section~\ref{sec:refseq} evaluates $S$ for the reference sequence
actually used here.

\subsection*{Inverting the transform}
The inverse transform, \texttt{dna\_inverse\_transform}, recovers
$p_i$ by rearranging Equation~\ref{eq:mod4}:
$p_i = (c_i - r_{i \bmod L} - c_{i-1}) \bmod 4$. This looks like a
direct mirror of the forward loop, but one detail is not symmetric.
The feedback value available while decrypting is $c_{i-1}$, the
previous \emph{ciphertext} nucleotide, not $p_{i-1}$, the previous
\emph{plaintext} nucleotide the forward loop would have used at that
point. At position $i$, the ciphertext nucleotide $c_i$ is already
known, since it is the input being decoded, so $c_{i-1}$ is always
available by the time it is needed, even before $p_i$ itself has been
computed. That is why the transform can be inverted one nucleotide at
a time: Listing~\ref{lst:mod4chain-inverse} makes this explicit, with
\texttt{prev} set to \texttt{c}, the value just consumed from the
input, never to \texttt{p}, the value just recovered.

\begin{lstlisting}[caption={Chained modulo-4 transform: the inverse
loop, rearranging Equation~\ref{eq:mod4} to recover $p_i$.}, label={lst:mod4chain-inverse}]
out_vals = []
prev = 0
for i, c in enumerate(cipher_vals):
    r = ref_vals[i % rlen]
    p = (c - r - prev) % 4
    out_vals.append(p)
    prev = c  # feedback uses the ciphertext nucleotide,
              # always known at decrypt time
\end{lstlisting}

\section{The Role of the Reference Sequence}
\label{sec:refseq}
The reference sequence $r$ is a short, fixed sequence of nucleotides.
It is public and does not need to stay secret. This is different from
the AES key $K$, which must stay secret for the cipher to work.

In this thesis, the reference sequence is stored in the file
\texttt{data/reference\_sequence.txt}. It has $L=61$ nucleotides, stays
the same for every run of the experiment, and reads

\begin{center}
{\ttfamily\small ATCGGATCTAGCTAGCTAGCTTACGGAATCCGGTACCTAGGCTATCGATCGTAGCTAGCTA}
\end{center}

Its base composition is close to even: 15 A, 15 C, 15 G, and 16 T, so
each base accounts for between 24.59\% and 26.23\% of the sequence. The
sequence contains the pattern \texttt{AGCT} five times, as two separate
back-to-back runs: at positions 9, 13 and 17, giving
\texttt{...AGCTAGCTAGCT...}, and again at positions 52 and 56, giving
\texttt{...AGCTAGCT...} (positions are zero-based).

For the accumulated mask defined in Section~\ref{sec:mod4chain}, what
matters is the sum over one period. Under the mapping of
Fig.~\ref{fig:byte-to-dna}, this sequence gives $S = 91$, and
$91 \equiv 3 \pmod 4$. Since 3 generates the whole of
$\mathbb{Z}_4$ under addition, the accumulated mask $R$ has period
$4L = 244$ rather than $L = 61$. This is a property of the sequence
rather than of the transform: a reference sequence whose nucleotide
values happen to sum to a multiple of 4 would give an accumulated mask
that repeats every $L$ positions instead.
Section~\ref{sec:results-correction} returns to this point, because the
alternative sequence compared there falls into exactly that case.

In Equation~\ref{eq:mod4}, the reference sequence supplies the value
$r_{i \bmod L}$ that gets added at position $i$. Most input files are
longer than $L=61$ nucleotides. Because of this, the index
$i \bmod L$ wraps around and reuses the reference sequence many times
over the course of one file.

\textbf{Example.} Take the first 5 nucleotides of the reference
sequence, \texttt{r = ATCGG} ($L=5$ for this example). Using the
mapping from Fig.~\ref{fig:byte-to-dna} ($A=0$, $T=1$, $C=2$, $G=3$),
this gives $r = [0, 1, 2, 3, 3]$. For an input stream of 12 nucleotides
($i = 0, \ldots, 11$), the index $i \bmod L$ cycles through
$0,1,2,3,4,0,1,2,3,4,0,1$. Fig.~\ref{fig:refseq-tiling} shows this
cycle. The same 5 values are reused more than twice across the 12 input
positions.

On its own, this reuse would repeat every $L$ positions and leave a
visible period in the output. Section~\ref{sec:mod4chain} explains how
the feedback term $c_{i-1}$ removes that periodicity.

\begin{figure}[htbp]
\centering
\begin{tikzpicture}[
  refbox/.style={draw, minimum width=0.9cm, minimum height=0.8cm, font=\ttfamily\small, fill=gray!12},
  tilebox/.style={draw, minimum width=0.65cm, minimum height=0.75cm, font=\ttfamily\small},
  >=Stealth
]
  % the reference sequence itself, L=5
  \foreach \i/\x/\v in {0/0/A,1/0.95/T,2/1.9/C,3/2.85/G,4/3.8/G} {
    \node[refbox] (r\i) at (\x,4.4) {\v};
    \node[font=\scriptsize, above=0.05cm of r\i] {$r_{\i}$};
  }
  \node[left=0.4cm of r0, font=\small, align=right] {Reference $r$\\($L{=}5$):};

  \draw[->] (2.4,3.9) -- node[right, font=\scriptsize, align=left] {tiled cyclically,\\wraps every $L$ positions} (2.4,3.0);

  % tiled row: i mod L for i=0..11
  \foreach \i/\x/\v in {0/0/A,1/0.65/T,2/1.3/C,3/1.95/G,4/2.6/G,5/3.25/A,6/3.9/T,7/4.55/C,8/5.2/G,9/5.85/G,10/6.5/A,11/7.15/T} {
    \node[tilebox] (t\i) at (\x,2.0) {\v};
    \node[font=\tiny, below=0.08cm of t\i] {$i{=}\i$};
  }
  \node[left=0.4cm of t0, font=\small, align=right] {Value used\\at position $i$:};

  % braces to mark the three cycles
  \draw[decorate, decoration={brace, amplitude=3pt}] (t0.north west) -- (t4.north east) node[midway, above=0.15cm, font=\scriptsize] {cycle 1};
  \draw[decorate, decoration={brace, amplitude=3pt}] (t5.north west) -- (t9.north east) node[midway, above=0.15cm, font=\scriptsize] {cycle 2};
  \draw[decorate, decoration={brace, amplitude=3pt}] (t10.north west) -- (t11.north east) node[midway, above=0.15cm, font=\scriptsize] {cycle 3};
\end{tikzpicture}
\caption{Tiling of a short, 5-nucleotide reference sequence
(\texttt{ATCGG}, the first 5 nucleotides of the reference sequence
used in this thesis) across 12 input positions. The index
$i \bmod L$ wraps around after every $L=5$ positions, so the same 5
values are reused for positions 5--9 and again, partially, for
positions 10--11. Without the feedback term in
Equation~\ref{eq:mod4}, this reuse would repeat every $L$ positions.}
\label{fig:refseq-tiling}
\end{figure}

\section{AES-GCM Fundamentals}
\label{sec:aesgcm}
AES-GCM is an authenticated cipher. It gives both confidentiality and
integrity at the same time. It does this by combining two mechanisms
that share the same block-cipher key $K$. CTR-mode encryption provides
confidentiality. GHASH, a chained multiply-accumulate operation over a
Galois field, provides integrity~\cite{nist80038d}. Three values play a
key role: the \emph{nonce} (a single-use value that seeds the counter,
generated randomly on every call), the hash subkey $H = E_K(0^{128})$,
and the authentication \emph{tag}. Reusing a nonce under the same key is
a serious security problem, because it causes the same keystream to be
reused across two different messages~\cite{nist80038d}. This thesis uses
AES-256 in GCM mode throughout: the key in \texttt{data/aes\_key.bin} is
32 bytes, and \texttt{src/classical\_crypto.py} wraps AES-256-GCM.

\begin{figure}[H]
\centering
\resizebox{\textwidth}{!}{%
\begin{tikzpicture}[
  ctrbox/.style={draw, minimum width=1.3cm, minimum height=0.7cm, font=\scriptsize},
  ekbox/.style={draw, minimum width=1.2cm, minimum height=0.8cm, font=\small},
  cbox/.style={draw, minimum width=1.1cm, minimum height=0.7cm, font=\small, fill=gray!12},
  xorc/.style={draw, circle, minimum size=0.6cm, inner sep=0pt, font=\Large},
  mulc/.style={draw, circle, minimum size=0.6cm, inner sep=0pt, font=\Large, fill=gray!12},
  >=Stealth
]
  % --- CTR-mode encryption row (top): three plaintext blocks ---
  \foreach \i/\x in {1/0,2/3,3/6} {
    \node[ctrbox] (ctr\i) at (\x,6.0) {$\mathrm{ctr}_0{+}\i$};
    \node[ekbox]  (ek\i)  at (\x,4.7) {$E_K$};
    \node[xorc]   (xr\i)  at (\x,3.2) {$\oplus$};
    \node[cbox]   (c\i)   at (\x,2.0) {$c_\i$};
    \draw[->] (ctr\i) -- (ek\i);
    \draw[->] (ek\i) -- (xr\i);
    \draw[->] (xr\i) -- (c\i);
    \node[left=0.15cm of xr\i, font=\small] {$m_\i$};
    \draw[->] ($(xr\i)+(-1.0,0)$) -- (xr\i);
  }

  % --- H generation inset (top right) ---
  \node[ekbox] (zek) at (9.2,5.4) {$E_K$};
  \node[font=\small, left=0.2cm of zek] {$0^{128}$};
  \draw[->] ($(zek)+(-1.0,0)$) -- (zek);
  \node[cbox, right=0.6cm of zek] (hgen) {$H$};
  \draw[->] (zek) -- (hgen);

  % --- Tag-mask counter (ctr0), off to the side of the CTR row ---
  \node[ctrbox] (ctr0) at (-2.6,6.0) {$\mathrm{ctr}_0$};
  \node[ekbox] (ek0) at (-2.6,4.7) {$E_K$};
  \draw[->] (ctr0) -- (ek0);

  % --- GHASH chain (bottom): multiply-accumulate with H ---
  \node[mulc] (m1) at (0,0.2) {$\otimes$};
  \node[xorc] (x1) at (1.5,0.2) {$\oplus$};
  \node[mulc] (m2) at (3,0.2) {$\otimes$};
  \node[xorc] (x2) at (4.5,0.2) {$\oplus$};
  \node[mulc] (m3) at (6,0.2) {$\otimes$};
  \node[xorc] (x3) at (7.5,0.2) {$\oplus$};
  \node[mulc] (m4) at (9,0.2) {$\otimes$};
  \node[xorc] (xt) at (10.5,0.2) {$\oplus$};
  \node[cbox] (tag) at (12.0,0.2) {$T$};
  \node[cbox] (len) at (7.5,2.0) {$L$};

  \draw[->] (c1.south) -- (m1.north);
  \draw[->] (m1) -- (x1);
  \draw[->] (c2.south) |- (x1.north);
  \draw[->] (x1) -- (m2);
  \draw[->] (m2) -- (x2);
  \draw[->] (c3.south) |- (x2.north);
  \draw[->] (x2) -- (m3);
  \draw[->] (m3) -- (x3);
  \draw[->] (len.south) -- (x3.north);
  \draw[->] (x3) -- (m4);
  \draw[->] (m4) -- (xt);
  \draw[->] (ek0.south) |- (xt.north);
  \draw[->] (xt) -- (tag);

  \foreach \n in {m1,m2,m3,m4} {
    \draw[->] ($(\n)+(0,-1.1)$) -- (\n) node[midway, right, font=\small] {$H$};
  }

  \node[font=\scriptsize, align=center, below=1.3cm of m1, xshift=0.4cm] {length block $L$\\encodes AAD/ciphertext lengths};
\end{tikzpicture}}
\caption{Structure of AES-GCM, based on the GCM
specification~\cite{nist80038d}. Top: CTR-mode encryption. Successive
counter values $\mathrm{ctr}_0{+}1, \mathrm{ctr}_0{+}2, \ldots$ are
encrypted with $E_K$ and XORed with plaintext blocks $m_i$. This
produces ciphertext blocks $c_i$. Top right: the hash subkey $H$ is
derived once, by encrypting an all-zero block. Bottom: GHASH processes
each ciphertext block and the length block $L$ in a chained
multiply-accumulate operation over $\mathrm{GF}(2^{128})$; each
$\otimes$ multiplies by $H$. The result is XORed with
$E_K(\mathrm{ctr}_0)$, the encryption of the counter value set aside
for the tag and never used for plaintext, to produce the authentication
tag $T$.}
\label{fig:aes-gcm}
\end{figure}

\section{Composing the Mechanisms into Two Pipelines}
\label{sec:pipelines}
The two mechanisms above are combined into two processing pipelines.
Both run on the same frozen corpus under the same AES key, so any
statistical difference between them can be attributed to the
DNA-transform step alone (Fig.~\ref{fig:pipelines}).

\begin{figure}[htbp]
\centering
\begin{tikzpicture}[
  node distance=0.7cm and 0.9cm,
  stage/.style={draw, minimum width=2.4cm, minimum height=1.1cm, align=center, font=\small, rounded corners=2pt},
  dnastage/.style={draw, minimum width=2.4cm, minimum height=1.1cm, align=center, font=\small, rounded corners=2pt, fill=gray!15},
  >=Stealth
]
  % Pipeline A (top)
  \node[stage] (a1) {Plaintext};
  \node[dnastage, right=of a1] (a2) {DNA Transform\\(mod-4 chain)};
  \node[stage, right=of a2] (a3) {AES-GCM\\Encryption};
  \node[stage, right=of a3] (a4) {Ciphertext};
  \node[font=\bfseries, above=0.2cm of a1, anchor=south west, xshift=-0.3cm] {Pipeline A};

  \draw[->] (a1) -- (a2);
  \draw[->] (a2) -- (a3);
  \draw[->] (a3) -- (a4);

  % Pipeline B (bottom), AES-GCM aligned under Pipeline A's AES-GCM stage
  \node[stage, below=2.0cm of a1] (b1) {Plaintext};
  \node[stage, right=of b1, xshift=2.9cm] (b2) {AES-GCM\\Encryption};
  \node[stage, right=of b2] (b4) {Ciphertext};
  \node[font=\bfseries, above=0.2cm of b1, anchor=south west, xshift=-0.3cm] {Pipeline B (control)};

  \draw[->] (b1) -- (b2);
  \draw[->] (b2) -- (b4);

  % measurement points (dotted verticals + labels)
  \draw[dotted] (a3.south) -- (b2.north);
  \node[font=\scriptsize, align=center] at ($(a3)!0.5!(b2)$) {post-encryption};

  \node[font=\scriptsize, align=center, above=0.3cm of a2] {post-DNA-transform\\(Pipeline A only)};
\end{tikzpicture}
\caption{The two processing pipelines under study. Pipeline A applies
the DNA-domain transform before AES-GCM encryption. Pipeline B is the
study control. It skips the DNA step and applies AES-GCM
alone to the same plaintext under the same key. Measurements (entropy,
chi-square, compression ratio, avalanche, correlation) are taken at the
input and after each pipeline's final stage. Pipeline A also has an
extra measurement point right after the DNA transform. Pipeline B has
no equivalent point.}
\label{fig:pipelines}
\end{figure}

% ============================================================
% CHAPTER 4 --- RELATED WORK
% ============================================================
\chapter{Related Work}
\label{ch:related}

\section{DNA-LWCS as the Architectural Inspiration}
Aqeel et al.~\cite{aqeel2025} propose DNA-LWCS. It combines keys
derived from publicly available DNA sequences, a
substitution-and-transposition stage, and ECC encryption (curve
secp256r1/P-256). Three keys are derived from a large public set of DNA
sequences. The substitution step applies XOR, modular addition, or AND,
picked at random per run. The transposition step then rotates the
result. ECC encrypts the message directly as elliptic-curve point
arithmetic: $C_1 = kP$, $C_2 = M + kQ$. The scheme is tested on three
sample images and compared against three other published
schemes~\cite{aqeel2025}.

Barman and Saha~\cite{dnaecc} follow a similar general idea: they encode
data as DNA sequences, then apply ECC-based encryption. Their motivation
is the large number of possible DNA sequences and their apparent
randomness. This thesis has not reviewed~\cite{dnaecc} at the same level
of detail as DNA-LWCS. Because of that, the detailed architectural
comparison below is built against DNA-LWCS. Where the comparison applies
to the general hybrid DNA-plus-ECC approach, it applies to both works.

\section{Divergences from This Thesis}
Table~\ref{tab:comparison} compares the two designs directly. The
classical layer is the clearest difference: DNA-LWCS uses ECC, an
asymmetric cipher, while this thesis uses AES-GCM, a symmetric cipher.
Section~\ref{sec:ecc-justification} explains why. The test data also
differs. DNA-LWCS reports its results on three sample images, while
this thesis uses text, JSON, and random data. The two works also test
different things. DNA-LWCS compares its whole system against other
published schemes. This thesis instead compares the same cipher with
and without the DNA-encoding step, which isolates the DNA step's own
effect.

\begin{table}[h]
\centering
\caption{Comparison of DNA-LWCS's architecture and this thesis's}
\label{tab:comparison}
\begin{tabular}{>{\raggedright\arraybackslash}p{3.2cm} >{\raggedright\arraybackslash}p{5.3cm} >{\raggedright\arraybackslash}p{5.3cm}}
\toprule
 & DNA-LWCS & This thesis \\
\midrule
Classical layer & ECC (asymmetric) & AES-GCM (symmetric) \\
Test data & 3 images & text/JSON/random \\
Baseline & other published schemes & same cipher, DNA on/off \\
Metrics & correlation, time, energy & entropy, chi-square, avalanche, correlation \\
Role of the DNA sequence & source of randomness for key generation & fixed transform parameter \\
Test corpus & data available on request, not fixed & frozen, checksummed corpus \\
\bottomrule
\end{tabular}
\end{table}

The reference DNA sequence plays a different role in each design. In
DNA-LWCS, DNA sequences are a source of randomness used to generate
keys. In this thesis, the reference sequence is a fixed, public
parameter mixed directly with the plaintext (Section~\ref{sec:refseq}).
The two test corpora also differ in how reproducible they are.
DNA-LWCS's data is described as available on request, with no fixed set
of files. This thesis uses a frozen, checksummed corpus, so the same
input files can be reused in later work.

\section{Justification for Replacing ECC with AES-GCM}
\label{sec:ecc-justification}
This thesis uses AES-GCM instead of ECC. 

The research question in Chapter~\ref{ch:aim} does not depend on a
specific classical cipher. It asks whether the DNA-encoding step
changes the data's statistical signature, and whether that change
survives classical encryption in general. Any classical cipher can
answer this question.

Pipeline A and
Pipeline B (Section~\ref{sec:pipelines} and
Fig.~\ref{fig:pipelines}) use the same key and the same cipher, with the
DNA step present or absent. This isolates the DNA transform's own
effect. AES-GCM is symmetric, so both pipelines can share one key
directly. ECC is asymmetric, and does not fit this same-key,
step-on/step-off comparison as directly.

DNA-LWCS's own use of ECC also takes a different path than the more
common pattern for bulk data. It encrypts the message directly through
point arithmetic ($C_2 = M + kQ$). In many practical systems, ECC is
used to exchange a key, and a symmetric cipher then encrypts the actual
message. This hybrid pattern, often called ECIES, is described in the
Standards for Efficient Cryptography~\cite{sec1}. Raw elliptic-curve
point operations scale less directly to arbitrary-length messages than
a block cipher does. Because of this, using AES-GCM here fits this
thesis's ablation design well, alongside DNA-LWCS's own, differently
motivated choice of ECC.

DNA-LWCS is this thesis's architectural inspiration, but the two designs
diverge from that starting point. DNA-LWCS asks a broad
question: is the whole combined system efficient and secure compared to
other whole systems. This thesis asks a narrower question: does the
DNA-encoding step itself leave a statistical fingerprint, isolated from
the cipher around it. DNA-LWCS's own paper does not test this narrower
question, since it never runs its system with the DNA step removed.
Running the original DNA-plus-ECC combination as a sanity check against
the AES-GCM baseline used here would not change what the ablation design
in Fig.~\ref{fig:pipelines} measures; Chapter~\ref{ch:limitations}
records this as a scope boundary.

\section{Differences in Evaluation Metrics}
DNA-LWCS and this thesis report different metrics, which is why a
direct, numeric comparison between the two is not possible. DNA-LWCS
focuses on correlation, encryption time, and energy use: metrics well
suited to comparing whole encryption systems against each other on
image data. This thesis instead focuses on entropy, chi-square, and the
avalanche effect: metrics suited to isolating one processing step's statistical contribution. DNA-LWCS also frames a higher correlation
value as a desirable outcome, tied to preserving data recovery. This
thesis instead measures correlation between plaintext and a later
stage's output as a check on how much linear structure survives that
stage, where a value near zero is the expected result once encryption
has been applied. These reflect two different evaluation goals, image-system
comparison versus step-by-step statistical isolation.

% ============================================================
% CHAPTER 5 --- RESEARCH METHODOLOGY
% ============================================================
\chapter{Research Methodology}
\label{ch:methodology}

\section{Threat Model}
\label{sec:threat-model}

This thesis measures statistical properties of ciphertext under a
specific adversary model. Chapter~\ref{ch:limitations}
notes that a statistical test is not a cryptanalytic security
proof. This section defines the model: what the attacker observes,
what they know, and what falls outside it.

\textbf{What the attacker sees.} Only the final transmitted output is
ever exposed: the nonce, the AES-GCM ciphertext body, and the
authentication tag (Section~\ref{sec:aesgcm}), concatenated as
\texttt{full\_aead\_bytes} (Section~\ref{sec:measurement-design}). The
\texttt{post\_dna} stage is an internal, in-memory value produced
immediately before AES-GCM encryption, used to isolate the DNA
transform's own behaviour for this thesis's ablation design
(Section~\ref{sec:pipeline-design}). This model assumes that value is
never transmitted and never leaves the encrypting device; that is an
assumption about how the pipeline is deployed, not a property the
construction itself guarantees.

\textbf{What the attacker knows.} Following Kerckhoffs's principle,
the attacker is assumed to know the full system design: the
DNA-domain transform (Section~\ref{sec:mod4chain}) and the reference
sequence $r$, which is explicitly public and requires no protection
(Section~\ref{sec:refseq}). The only value the attacker does not know
is the AES key $K$, generated once and kept secret for the duration
of the experiment.

\textbf{What the attacker is trying to do.} This thesis considers a
passive, ciphertext-only adversary: one who observes transmitted
ciphertexts and attempts to learn something about the corresponding
plaintext, for example, distinguishing structured text or JSON
from random data without submitting chosen plaintexts or
interacting with an encryption oracle. This matches what
Chapter~\ref{ch:results}'s metrics test: whether \texttt{post\_classical}
output alone reveals a statistical difference across data classes. It
does not model active attacks like forgery, tampering, or
chosen-ciphertext manipulation, which AES-GCM's authentication tag is designed to
prevent~\cite{nist80038d} and which this thesis does not evaluate.

\textbf{Out of scope.} Side-channel attacks (timing, power) against either the DNA
transform or the AES-GCM implementation are out of scope, as are
key-recovery attacks against AES itself and attacks on key
distribution or management.

The attacker in this model never sees \texttt{post\_dna} output, only
\texttt{post\_classical}. Measurements taken at \texttt{post\_dna} are
therefore diagnostics of the transform, not statements about what this
adversary can observe.

\section{Pipeline Comparison Design}
\label{sec:pipeline-design}
This thesis compares two pipelines directly. Pipeline A applies the DNA-domain transform,
then AES-GCM. Pipeline B applies AES-GCM alone. Both pipelines use the
same fixed AES key and the same input file. The only difference between
them is the DNA-transform step. The design goal is to isolate the DNA
transform's own effect on the data.

Each pipeline exposes its intermediate stages. Pipeline A exposes three stages:
\texttt{input}, \texttt{post\_dna}, and \texttt{post\_classical}.
Pipeline B exposes two stages: \texttt{input} and
\texttt{post\_classical}. Every metric in Section~\ref{sec:metrics} is
computed at each available stage. This gives three measurement points
for Pipeline A and two for Pipeline B, matching Fig.~\ref{fig:pipelines}.

The two pipelines are deterministic, except for one part. AES-GCM
generates a fresh random nonce on every call
(Section~\ref{sec:aesgcm}). Because of this, \texttt{post\_classical}
output differs between runs, even for the same input and key. The
\texttt{post\_dna} stage has no randomness at all: the same input and
reference sequence always produce the same output.
Section~\ref{sec:measurement-design} explains why this difference matters for one
of the metrics.

\section{Test Data Corpus}
\label{sec:corpus}
The experiment uses a frozen corpus of six files, covering three data classes, each at two sizes, small (about 10 KB)
and large (about 200 KB).

\textbf{Text.} English prose, a slice taken from the opening of Jane
Austen's \emph{Pride and Prejudice} (Project Gutenberg eBook \#1342,
public domain).

\textbf{JSON.} Synthetic structured log records, generated with two
different random seeds, so the large file is not a repeat of the
small one.

\textbf{Random.} Bytes from \texttt{os.urandom}, used as a
high-entropy control.

Table~\ref{tab:corpus} lists the six files with their size and Shannon
entropy (Section~\ref{sec:metrics}).

\begin{table}[h]
\centering
\caption{The frozen test corpus}
\label{tab:corpus}
\begin{tabular}{lrr}
\toprule
File & Size (bytes) & Entropy (bits/byte) \\
\midrule
text\_small.txt & 10238 & 4.7268 \\
text\_large.txt & 204799 & 4.6162 \\
json\_small.json & 10313 & 4.5202 \\
json\_large.json & 205012 & 4.5324 \\
random\_small.bin & 10240 & 7.9804 \\
random\_large.bin & 204800 & 7.9991 \\
\bottomrule
\end{tabular}
\end{table}

Text and JSON land at 4.52 to 4.73 bits/byte. Random data lands at 7.98
to 8.00 bits/byte. This gap gives the experiment two clearly separated
entropy classes to compare, before any transform is applied.

The corpus is frozen and checksummed.
\texttt{data/corpus/CHECKSUMS.sha256} records the SHA-256 hash of each
file, so the exact byte content can be checked again later with
\texttt{shasum -a 256 -c data/corpus/CHECKSUMS.sha256}.
\texttt{build\_corpus.py} downloads its source text live from the
internet, so without a checksum, a later change to the Gutenberg file
could silently change the corpus that later results are compared
against.

\section{Metrics Used}
\label{sec:metrics}
Five metrics are computed at each pipeline stage: entropy, chi-square,
compression ratio, correlation, and avalanche effect. Two more,
throughput and size overhead, come from the benchmark harness
(\texttt{src/benchmark.py}), which runs each pipeline end to end under
timing instrumentation rather than at a single stage.

\textbf{Shannon entropy.} Bits per byte, computed from a 256-bin
histogram of byte values:
\begin{equation}
H = -\sum_{i=0}^{255} p_i \log_2 p_i
\label{eq:entropy}
\end{equation}
where $p_i$ is the observed frequency of byte value $i$. Higher values
mean the byte distribution is closer to uniform (bins with a zero
count are excluded before summing, since $\log_2 0$ is undefined).

\textbf{Chi-square statistic.} Tests whether the byte-value
distribution matches a uniform distribution over all 256 values, via
\texttt{scipy.stats.chisquare}. A larger statistic means a larger
departure from uniformity. Because the raw statistic scales with
message length, comparisons in Chapter~\ref{ch:results} are also
reported as the standardised deviate
$(\chi^2 - \nu)/\sqrt{2\nu}$ with $\nu = 255$ degrees of freedom, which
expresses the departure in units of the statistic's own standard
deviation and is comparable across files of similar size.

\textbf{Compression ratio.} The ratio of gzip-compressed size to
original size. Lower values mean more redundancy. Values near or above
1.0 mean the data is close to incompressible.

\textbf{Pearson correlation.} The correlation coefficient between the
byte values of two equal-length byte strings, usually the original
input and a later stage's output. Values near zero mean little linear
relationship remains between the two.

\textbf{Avalanche effect.} One random bit is flipped in a copy of the
input. Both the original and the flipped copy are run through the same
stage. The Hamming distance between the two outputs, as a percentage of
total bits, is recorded, and the mean across trials is the reported
avalanche percentage. Two trial counts are used in this thesis: the
full-pipeline runs recorded in \texttt{results/experiment\_results.csv}
use the harness default of 25 trials, while the \texttt{post\_dna}
measurement reported in Section~\ref{sec:results-avalanche} uses 2,000
trials, since that is the measurement the analysis actually rests on.
This one-bit-flip method follows the same design used to test lightweight
block ciphers~\cite{inru}. A textbook-ideal avalanche effect lands
close to 50\%. Listing~\ref{lst:avalanche-algo} gives the trial loop
(\texttt{src/avalanche.py}): flip one random bit, run both copies
through \texttt{pipeline\_encrypt\_fn}, and score the Hamming distance
between the two outputs as a percentage of total output bits.

\begin{lstlisting}[caption={Avalanche-effect trial loop: one-bit flip,
run both copies through the pipeline, score the Hamming distance
between the outputs (\texttt{src/avalanche.py}).},
label={lst:avalanche-algo}]
for _ in range(n_trials):
    flipped = bytearray(data)
    bit_index = random.randrange(len(data) * 8)
    byte_index, bit_in_byte = divmod(bit_index, 8)
    flipped[byte_index] ^= 1 << bit_in_byte
    flipped = bytes(flipped)

    ciphertext_original = pipeline_encrypt_fn(data)
    ciphertext_flipped = pipeline_encrypt_fn(flipped)

    n_bits = min(len(ciphertext_original), len(ciphertext_flipped)) * 8
    hamming_distance = _hamming_distance_bits(
        ciphertext_original[:n_bits // 8], ciphertext_flipped[:n_bits // 8]
    )
    percentages.append(100.0 * hamming_distance / n_bits)
\end{lstlisting}

\textbf{Throughput.} Data size in bytes divided by total pipeline time
in seconds, giving bytes per second. This is measured for every pipeline
run, but Chapter~\ref{ch:results} does not report it as a result:
Chapter~\ref{ch:limitations} explains why a raw timing number from this
thesis's specific, unoptimized implementation is not reported as a
finding about the DNA-domain approach in general.

\textbf{Size overhead.} The percentage increase in size from plaintext
to the encrypted output:
\begin{equation}
\text{overhead} = \frac{\text{len(encrypted)} - \text{len(plaintext)}}{\text{len(plaintext)}} \times 100
\label{eq:overhead}
\end{equation}
Section~\ref{sec:measurement-design} explains that which bytes count as "the
encrypted output" changes this number a great deal.

Timing uses \texttt{time.perf\_counter()}, with 20 repeated calls per
measurement and the first call discarded as a warm-up. The median and
interquartile range are reported, following the same repeated-measurement
approach used in lightweight-cipher performance surveys~\cite{ascon-survey}.

\section{Measurement Design: Avalanche and Size Overhead}
\label{sec:measurement-design}
Two of the metrics above depend not only on which pipeline stage
produced the bytes, but on exactly which bytes get passed in. How that
question is resolved for avalanche effect and size overhead shapes the
numbers Chapter~\ref{ch:results} reports for both.

\textbf{Avalanche effect: which stage to measure.} Section~\ref{sec:pipeline-design}
noted that \texttt{post\_classical} output is randomized by AES-GCM's
per-call nonce, while \texttt{post\_dna} is fully deterministic. This
matters directly for avalanche effect, since the metric is defined as
the Hamming distance between two outputs that differ by one flipped
input bit. AES-GCM's nonce guarantees that any two independently
encrypted ciphertexts look statistically unrelated regardless of the
plaintext (Section~\ref{sec:aesgcm}), so measuring avalanche on
Pipeline A's \texttt{post\_classical} ciphertext would report close to
50\% whether or not the input bit was actually flipped: the nonce
alone accounts for nearly all of the observed difference. Measuring at
\texttt{post\_dna} instead isolates the quantity this thesis actually
wants: with no nonce in the way, the Hamming distance between two
\texttt{post\_dna} outputs reflects only the one flipped bit and the
transform's own propagation behaviour.
Section~\ref{sec:results-avalanche} reports that measurement and
derives its expected value from Equation~\ref{eq:prefixsum}.

\textbf{Size overhead: which bytes count as output.} Equation~\ref{eq:overhead}
depends on \texttt{len(encrypted)}, and AES-GCM's output is really
three distinct pieces: the ciphertext body, the nonce, and the
authentication tag (Section~\ref{sec:aesgcm}). AES-GCM's CTR-mode
encryption keeps the ciphertext body exactly the same length as the
plaintext, so measuring overhead against the ciphertext body alone
returns 0.0\% for every input, independent of the DNA transform or
anything else about the data.

RFC 5116's own definition of AEAD ciphertext expansion counts only the
authentication tag: an AEAD\_AES\_256\_GCM ciphertext, the registered
algorithm matching the AES-256 key used here, is defined as exactly 16
octets longer than its plaintext, the tag alone (\cite{rfc5116},
Sections 5.1--5.2). The nonce is deliberately left out of that
definition. RFC 5116 treats it as something the application must get to
the receiver by whatever means it chooses, noting in Section 2.1 that
the nonce may be stored or transported with the ciphertext, or else
reconstructed by the decrypter. This project's own pipeline
(\texttt{src/pipeline.py}) takes the first option: it bundles
\texttt{nonce}, \texttt{ciphertext}, and \texttt{tag} together as the
one unit that gets stored or transmitted (Listing~\ref{lst:aead-overhead}).
Counting the nonce alongside the tag, then, reflects this project's own
transmitted-output convention, not a universal AEAD accounting rule.
Counted this way, overhead becomes size-dependent, as
Section~\ref{sec:results-overhead} reports: about 0.31\% for a 10 KB
file and about 0.016\% for a 200 KB file, since the nonce and tag add a
fixed number of bytes regardless of message size. The nonce length that
produces those figures departs from both cited standards;
Section~\ref{sec:limits-nonce} records that as a limitation.

\texttt{encrypt\_aes\_gcm} itself returns the three pieces separately,
as a tuple, not concatenated (Listing~\ref{lst:aead-overhead}, top).
The concatenation into a single byte string, the object Equation~\ref{eq:overhead} actually measures the length of, happens on
the caller's side, once a decision has been made about which bytes
count as "the encrypted output."

\begin{lstlisting}[caption={What gets concatenated into the full-AEAD
byte string that Equation~\ref{eq:overhead} is measured against
(\texttt{src/classical\_crypto.py},
\texttt{scripts/check\_aead\_overhead\_convention.py}).},
label={lst:aead-overhead}]
def encrypt_aes_gcm(data: bytes, key: bytes) -> tuple[bytes, bytes, bytes]:
    cipher = AES.new(key, AES.MODE_GCM)
    ciphertext, tag = cipher.encrypt_and_digest(data)
    return cipher.nonce, ciphertext, tag

nonce, ciphertext, tag = encrypt_aes_gcm(plaintext, key)
full_aead_bytes = nonce + ciphertext + tag  # what "encrypted output" means here
\end{lstlisting}

\section{Repeated-Sample Measurement for Entropy, Chi-Square, and Correlation}
\label{sec:repeated-sample}
Avalanche effect is measured with 2,000 trials and reported as a mean
with a standard deviation (Section~\ref{sec:results-avalanche}).
Entropy, chi-square, and correlation are not: each is computed once per
data class, on the one frozen small file for that class. A single
number carries no sense of how much it could have varied under a
different draw of the same kind of data, so this section adds the same
repeated-sample treatment already used for avalanche.

For each data class, ten independent samples of about 10 KB are
generated, separately from the frozen corpus (Section~\ref{sec:corpus})
so the frozen files stay untouched. Random samples are ten fresh
\texttt{os.urandom()} draws. JSON samples are ten freshly generated logs
from the same generator that built the frozen corpus
(\texttt{scripts/build\_corpus.py}), each with its own seed. Text
samples are ten non-overlapping slices spaced evenly across the frozen
\texttt{text\_large.txt}, so no two slices share content. Each sample is
run through Pipeline A, and entropy, chi-square, and correlation are
computed at every stage, using the same functions as
Section~\ref{sec:metrics}. \texttt{scripts/measure\_repeated\_sample\_stats.py}
implements this; \texttt{results/repeated\_sample\_stats.csv} holds the
full set of means and standard deviations, reported in
Section~\ref{sec:results-repeated-sample}.

Where Chapter~\ref{ch:results} compares two data classes at the same
stage, it does so with Welch's unequal-variance $t$-test over these ten
samples per class, reporting the difference in means divided by its
standard error rather than a bare ratio of the raw statistics.

\section{Computational Environment}
\label{sec:environment}
All measurements reported in Chapter~\ref{ch:results} were produced on
a single machine: Apple M2 (arm64) running macOS (Darwin 23.2.0). The
software stack is Python 3.11.7 in the project's virtual environment,
with PyCryptodome 3.23.0 for AES-GCM, NumPy 2.4.6, and SciPy 1.17.1 for
the entropy and chi-square computations. The timing methodology in
Section~\ref{sec:metrics} is sensitive to all of these; the statistical
results are not, since the pipelines are deterministic given the same
input, key, and reference sequence.
% TODO: add repository URL and the commit hash the reported results
% were generated from, once the project is archived.

% ============================================================
% CHAPTER 6 --- IMPLEMENTATION
% ============================================================
\chapter{Implementation}
\label{ch:implementation}

\section{Module Structure}
\label{sec:module-structure}
The implementation is organized into small, single-purpose modules
under \texttt{src/}, each with its own test file under \texttt{tests/}.

\begin{itemize}
\item \texttt{dna\_codec.py}: converts between bytes, bits, and
nucleotide sequences (Section~\ref{sec:bytes-to-dna}). Functions:
\texttt{bytes\_to\_bits}, \texttt{bits\_to\_dna}, \texttt{dna\_to\_bits},
\texttt{bits\_to\_bytes}.
\item \texttt{dna\_transform.py}: the chained modulo-4 transform
(Section~\ref{sec:mod4chain}). Functions: \texttt{dna\_transform},
\texttt{dna\_inverse\_transform}.
\item \texttt{classical\_crypto.py}: a thin wrapper around AES-256-GCM
(Section~\ref{sec:aesgcm}). Functions: \texttt{encrypt\_aes\_gcm},
\texttt{decrypt\_aes\_gcm}.
\item \texttt{pipeline.py}: composes the two pipelines
(Section~\ref{sec:pipeline-design}). Functions:
\texttt{pipeline\_a\_dna\_\allowbreak plus\_classical},
\texttt{pipeline\_b\_classical\_only}, and their matching decrypt
functions.
\item \texttt{metrics.py}: entropy, chi-square, compression ratio, and
correlation (Section~\ref{sec:metrics}).
\item \texttt{avalanche.py}: the avalanche-effect test.
\item \texttt{benchmark.py}: timing, throughput, and size-overhead
measurement.
\end{itemize}

\section{Revision History of the DNA Transform Design}
\label{sec:revision-history}
The DNA-domain transform went through three designs during implementation. Each design was kept or dropped based on what its output measured.

\textbf{Design 1: SHA-256 mask plus XOR.} A mask was derived from a DNA
sequence by hashing it, then combined with the data using XOR. Measured
on 4,000-byte samples, text and random data both converged to about
7.96 bits/byte of entropy after the transform. A cryptographically strong mask pushes every input toward the
same near-maximal entropy, no matter what the input originally looked
like. The experiment depends on measuring a difference between low-entropy and high-entropy input (Chapter~\ref{ch:aim}).

Design 1
erased that difference entirely, so it was rejected, even though every
round-trip test passed: applying the transform, then its inverse,
recovered the exact original bytes for every case tested, including
random-length inputs, empty input, every single possible byte value,
and both text and binary samples. Round-trip correctness on its own
says nothing about whether a design preserves the entropy-class
signal this thesis measures, which is why Design 1 could pass every
such test and still be rejected.

\textbf{Design 2: tiled mask, no feedback.} The reference sequence was
tiled directly across the data, without hashing and without any
feedback between positions. This preserved an entropy gap (text at
6.805 bits/byte versus random at 7.951 bits/byte), but introduced a
different problem. Ciphertext-only autocorrelation at a lag equal to
the reference sequence's length measured 0.835 on structured text,
versus about $-0.008$ for random data, a severe weakness with the same
structural cause as a classic repeating-key XOR cipher: the same short
mask value repeats predictably every $L$ positions, and this
repetition survives into the output. Design 2 was rejected for this
reason too.

\textbf{Design 3: chained modulo-4 addition (final).} Adding a feedback
term, $c_{i-1}$, to Equation~\ref{eq:mod4} means each output nucleotide
depends on the entire stream so far, not only its position modulo $L$.

Measured on the same 4,000-byte samples, autocorrelation at lag $L$
dropped to about 0.017 to 0.018, back to chance level. The entropy gap also survived: text measured 7.272 bits/byte against random's 7.953. Both results held, so this design was kept.

Two caveats attach to those Design 3 numbers. First, the
7.272-versus-7.953 gap was measured using a
64-nucleotide sample sequence hardcoded in a test script, not the frozen reference sequence that every pipeline run in this thesis uses (Section~\ref{sec:refseq}); Section~\ref{sec:results-correction}
reports the gap measured against the real reference sequence
and corpus, and explains why the two numbers differ so much. Second,
the autocorrelation figure comes from the same test-script setup, and
was not re-measured against the frozen reference sequence and corpus
once those were in place. Given how strongly the reference sequence
turns out to affect the entropy gap
(Section~\ref{sec:results-correction}), that re-measurement is listed
as future work in Section~\ref{sec:future-work}.

% ============================================================
% CHAPTER 7 --- RESULTS
% ============================================================
\chapter{Results}
\label{ch:results}


\section{Entropy, Chi-Square, and Compression Signal}
\label{sec:results-entropy}
Table~\ref{tab:results-small} reports entropy, chi-square, and
compression ratio for the three small corpus files, at each stage of
Pipeline A. The same pattern holds at the large file sizes.

\begin{table}[h]
\centering
\caption{Entropy, chi-square, and compression ratio by stage, Pipeline A, small files}
\label{tab:results-small}
\begin{tabular}{llrrr}
\toprule
File & Stage & Entropy (bits/byte) & Chi-square & Compression ratio \\
\midrule
text\_small.txt & input & 4.7268 & 139108.219 & 0.4301 \\
 & post\_dna & 7.9265 & 1119.518 & 0.9833 \\
 & post\_classical & 7.9805 & 273.103 & 1.0022 \\
json\_small.json & input & 4.5202 & 179852.453 & 0.1817 \\
 & post\_dna & 7.8786 & 1960.032 & 0.8459 \\
 & post\_classical & 7.9814 & 267.597 & 1.0022 \\
random\_small.bin & input & 7.9804 & 271.800 & 1.0022 \\
 & post\_dna & 7.9834 & 234.250 & 1.0022 \\
 & post\_classical & 7.9827 & 242.300 & 1.0022 \\
\bottomrule
\end{tabular}
\end{table}

At input, text and JSON measure far below random on all three metrics,
as expected from the corpus design itself (Section~\ref{sec:corpus}).
Fig.~\ref{fig:byte-histogram} shows why, at the level of individual
byte values. Text's raw bytes
cluster in a few narrow bins, mostly the printable ASCII range. Random
bytes spread almost evenly across all 256 possible values from the
start.

\begin{figure}[H]
\centering
\begin{tikzpicture}[
  every axis/.style={
    width=5.3cm, height=4.2cm,
    ybar, bar width=3.5pt,
    ymin=0, ymax=0.5,
    xtick=\empty,
    ytick={0,0.1,0.2,0.3,0.4,0.5},
    yticklabel style={font=\scriptsize},
    title style={font=\scriptsize},
    axis line style={gray!60},
  }
]
\begin{axis}[name=hist1, title={Input: text\_small.txt},
  ylabel={fraction of bytes}, y label style={font=\scriptsize}]
\addplot+[fill=seqblue!70,draw=none] coordinates {
(1,0.0521)(2,0.0)(3,0.1854)(4,0.006)(5,0.0223)(6,0.0096)(7,0.4611)(8,0.222)(9,0.0139)(10,0.0139)(11,0.0)(12,0.0)(13,0.0)(14,0.0)(15,0.0139)(16,0.0)
};
\end{axis}
\begin{axis}[name=hist2, at={(hist1.right of south east)}, anchor=left of south west,
  title={Post-DNA: text\_small.txt}, yticklabel=\empty]
\addplot+[fill=seqblue!70,draw=none] coordinates {
(1,0.0618)(2,0.0583)(3,0.0611)(4,0.0651)(5,0.0669)(6,0.0621)(7,0.0614)(8,0.0635)(9,0.0644)(10,0.0633)(11,0.066)(12,0.0616)(13,0.0612)(14,0.0588)(15,0.0621)(16,0.0622)
};
\end{axis}
\begin{axis}[name=hist3, at={(hist2.right of south east)}, anchor=left of south west,
  title={Input: random\_small.bin}, yticklabel=\empty]
\addplot+[fill=seqblue!70,draw=none] coordinates {
(1,0.0648)(2,0.0646)(3,0.0595)(4,0.0586)(5,0.058)(6,0.0593)(7,0.0604)(8,0.066)(9,0.062)(10,0.0635)(11,0.0676)(12,0.0606)(13,0.0646)(14,0.062)(15,0.0658)(16,0.0626)
};
\end{axis}
\end{tikzpicture}
\caption{Byte-value histograms (16 bins, 16 consecutive byte values
each), computed directly from the corpus files and from Pipeline
A's actual \texttt{post\_dna} output. The coarser binning here is for
legibility only; every chi-square statistic in this chapter uses the
full 256-bin histogram of Section~\ref{sec:metrics}. Left:
\texttt{text\_small.txt}'s
raw input bytes, concentrated in a few bins. Middle: the same file
after the DNA transform, spread almost evenly across all 16 bins.
Right: \texttt{random\_small.bin}'s raw input bytes, already close to
even. The transform visibly flattens text's distribution, though
Fig.~\ref{fig:chisquare-heatmap} shows a small difference from random
survives underneath this flattened shape.}
\label{fig:byte-histogram}
\end{figure}

At \texttt{post\_dna}, entropy for text and JSON rises close to
random's level (7.9265 and 7.8786 bits/byte, against random's 7.9834),
but chi-square tells a sharper story: text still measures 1119.518,
about 4.8 times random's 234.250, and JSON measures 1960.032, about 8.4
times random's value. With 255 degrees of freedom (256 byte values
minus one), the chi-square critical value at $\alpha=0.05$ is 293.25:
text's and JSON's \texttt{post\_dna} statistics both clear it by a wide
margin, while every \texttt{post\_classical} value and every
\texttt{random\_small.bin} value at any stage falls below it. Expressed
as the standardised deviate defined in Section~\ref{sec:metrics}, the
same three \texttt{post\_dna} values are $38.3$ for text and $75.5$ for
JSON, against $-0.9$ for random data: text and JSON depart from
uniformity by tens of standard deviations, while random data is
indistinguishable from uniform. Fig.~\ref{fig:chisquare-heatmap} shows this
pattern across all three stages at once.

\begin{figure}[H]
\centering
\begin{tikzpicture}[
  cell/.style={minimum width=2.7cm, minimum height=1.3cm, draw=black!40, font=\small},
]
  \node[font=\small\bfseries] at (2.7,1.0) {input};
  \node[font=\small\bfseries] at (5.4,1.0) {post\_dna};
  \node[font=\small\bfseries] at (8.1,1.0) {post\_classical};
  \node[font=\small\bfseries, anchor=east] at (-0.2,0.0) {text\_small};
  \node[font=\small\bfseries, anchor=east] at (-0.2,-1.4) {json\_small};
  \node[font=\small\bfseries, anchor=east] at (-0.2,-2.8) {random\_small};

  \node[cell, fill=seqblue!92] at (2.7,0.0) {\textcolor{white}{139108.2}};
  \node[cell, fill=seqblue!34] at (5.4,0.0) {1119.5};
  \node[cell, fill=seqblue!17] at (8.1,0.0) {273.1};

  \node[cell, fill=seqblue!95] at (2.7,-1.4) {\textcolor{white}{179852.5}};
  \node[cell, fill=seqblue!41] at (5.4,-1.4) {1960.0};
  \node[cell, fill=seqblue!17] at (8.1,-1.4) {267.6};

  \node[cell, fill=seqblue!17] at (2.7,-2.8) {271.8};
  \node[cell, fill=seqblue!15] at (5.4,-2.8) {234.3};
  \node[cell, fill=seqblue!15] at (8.1,-2.8) {242.3};
\end{tikzpicture}
\caption{Chi-square statistic (Pipeline A) across the three small
corpus files and pipeline stages, from Table~\ref{tab:results-small}.
Darker cells mean a larger departure from a uniform byte distribution;
shading is scaled by $\log_{10}$ of the statistic, since the raw values
span more than three orders of magnitude, but the number shown in each
cell is untransformed statistic. The DNA transform
(\texttt{post\_dna} column) drops text's and JSON's chi-square by more
than two orders of magnitude from their input values, but both still
measure clearly higher than random data's. That remaining gap
disappears once AES-GCM (\texttt{post\_classical} column) runs: all
nine cells in that column land within a narrow band, regardless of
data class.}
\label{fig:chisquare-heatmap}
\end{figure}

Compression ratio follows the same direction as chi-square. Text goes
from 0.4301 (input) to 0.9833 (post\_dna), close to random's 1.0022.
JSON retains more compressibility than text at the same stage: 0.1817
(input) to 0.8459 (post\_dna), still visibly below random's ratio. At
\texttt{post\_classical}, entropy, chi-square, and compression ratio
all land within a narrow band across all three data classes and both
pipelines. \texttt{results/experiment\_results.csv} confirms Pipeline
B's \texttt{post\_classical} numbers (no DNA stage at all) fall in the
same band as Pipeline A's, for every file. Whatever signal survives the
DNA transform alone is gone by the time AES-GCM has run.

\section{Correlation with the Original Input}
\label{sec:results-correlation}
Section~\ref{sec:metrics} defined Pearson correlation as a check on
how much linear structure between a stage's output and the original
input survives that stage, with values near zero meaning little
structure remains. Table~\ref{tab:results-correlation} reports it for
the three small corpus files. At \texttt{input}, correlation is
trivially 1.0, since a stage is being compared against itself. At both
\texttt{post\_dna} and \texttt{post\_classical}, every value sits
within about $\pm 0.02$ of zero, consistent across all three data
classes: essentially no linear structure survives either stage.

\begin{table}[h]
\centering
\caption{Pearson correlation with the original input, by stage, Pipeline A, small files}
\label{tab:results-correlation}
\begin{tabular}{lrr}
\toprule
File & Stage & Correlation \\
\midrule
text\_small.txt & input & 1.0000 \\
 & post\_dna & -0.0024 \\
 & post\_classical & 0.0096 \\
json\_small.json & input & 1.0000 \\
 & post\_dna & -0.0070 \\
 & post\_classical & 0.0009 \\
random\_small.bin & input & 1.0000 \\
 & post\_dna & -0.0051 \\
 & post\_classical & -0.0188 \\
\bottomrule
\end{tabular}
\end{table}

Correlation does not distinguish the DNA transform's own effect from
AES-GCM's: both stages already look uncorrelated with the input, so
this metric does not separately confirm the entropy-class signal
chi-square finds at \texttt{post\_dna}. It rules out one specific kind
of structure, a linear relationship between input and output bytes,
without being sensitive enough to detect the smaller, non-linear signal
chi-square does find.

\section{Repeated-Sample Results}
\label{sec:results-repeated-sample}
Table~\ref{tab:results-repeated-sample} reports the ten-sample mean and
standard deviation for entropy, chi-square, and correlation at each
stage of Pipeline A, from the repeated-sample measurement described in
Section~\ref{sec:repeated-sample}.

\begin{table}[h]
\centering
\caption{Entropy, chi-square, and correlation, mean $\pm$ stdev over 10 independent $\sim$10 KB samples per data class, Pipeline A}
\label{tab:results-repeated-sample}
\begin{tabular}{llrrr}
\toprule
Data class & Stage & Entropy (bits/byte) & Chi-square & Correlation \\
\midrule
text & input & $4.6079 \pm 0.0593$ & $149039.7 \pm 5595.8$ & $1.0000 \pm 0.0000$ \\
 & post\_dna & $7.9239 \pm 0.0040$ & $1138.5 \pm 54.9$ & $-0.0047 \pm 0.0069$ \\
 & post\_classical & $7.9826 \pm 0.0007$ & $246.3 \pm 10.0$ & $0.0027 \pm 0.0077$ \\
json & input & $4.5218 \pm 0.0034$ & $179696.3 \pm 171.4$ & $1.0000 \pm 0.0000$ \\
 & post\_dna & $7.8908 \pm 0.0063$ & $1774.8 \pm 101.0$ & $0.0040 \pm 0.0094$ \\
 & post\_classical & $7.9830 \pm 0.0012$ & $242.3 \pm 18.2$ & $-0.0038 \pm 0.0089$ \\
random & input & $7.9819 \pm 0.0015$ & $257.1 \pm 21.7$ & $1.0000 \pm 0.0000$ \\
 & post\_dna & $7.9827 \pm 0.0015$ & $244.6 \pm 21.1$ & $0.0033 \pm 0.0105$ \\
 & post\_classical & $7.9815 \pm 0.0019$ & $260.9 \pm 25.4$ & $0.0003 \pm 0.0096$ \\
\bottomrule
\end{tabular}
\end{table}

The single-sample numbers already reported in Table~\ref{tab:results-small}
and Table~\ref{tab:results-correlation} sit close to these means: text's
\texttt{post\_dna} chi-square was 1119.518 against a ten-sample mean of
1138.5, JSON's was 1960.032 against 1774.8, and random's \texttt{input}
chi-square was 271.800 against 257.1. All three fall within about one
standard deviation of the ten-sample mean, so the single frozen file was
not an unusual draw.

The gap between data classes is larger than the spread within a class.
Applying Welch's unequal-variance $t$-test to the \texttt{post\_dna}
chi-square samples, text's mean of 1138.5 exceeds random's 244.6 by
48.1 standard errors, and JSON's mean of 1774.8 exceeds it by 46.9
standard errors ($n=10$ per class in both cases). Differences of that
size are far outside anything the observed within-class spread could
produce, which is the evidence H1 rests on.
Correlation stays within about
$\pm 0.01$ of zero at \texttt{post\_dna} and \texttt{post\_classical} for
every data class, in both mean and standard deviation, matching the
single-sample finding in Section~\ref{sec:results-correlation} and
confirming it holds across repeated draws, not only for the one frozen
file.

\section{Correction of the Original Entropy-Gap Estimate}
\label{sec:results-correction}
Section~\ref{sec:revision-history} reported an entropy gap of 7.272
bits/byte (text) versus 7.953 bits/byte (random) for the final,
chained modulo-4 design. That measurement used a 64-nucleotide sample
sequence hardcoded in a test script, not the real, frozen 61-nucleotide
reference sequence that every pipeline run in this thesis actually uses
(Section~\ref{sec:refseq}). Table~\ref{tab:results-small} shows the
real gap is much smaller: 7.9265 bits/byte (text, \texttt{post\_dna})
versus 7.9834 bits/byte (random, \texttt{post\_dna}), the same
\texttt{post\_dna}-to-\texttt{post\_dna} comparison Section~\ref{sec:results-entropy}
uses, a difference of only 0.0569 bits/byte.

This difference was traced to the reference sequence itself, not to a
bug in the transform or the corpus, by swapping reference sequence and
input data against each other. Holding the real corpus fixed and
substituting the old 64-nucleotide sample reference reproduces a 0.22
bits/byte gap, close to the original, larger estimate. Truncating the
real corpus down to the same 4,000-byte sample size, while keeping the
real reference, isolates the effect of sample size alone: about a 0.05
bit gap, so sample size accounts for very little of the difference next
to the roughly 0.17 bit swing the reference-sequence swap produces.
Substituting the real reference into the older hand-built samples gave
an intermediate and noticeably more variable gap, but that particular
comparison was run as an ad hoc analysis and was never persisted as a
script, so no reproducible figure is quoted for it here; re-running all
of these swaps under a committed script is listed in
Section~\ref{sec:future-work}. On the two swaps that are reproducible,
the reference sequence, not the corpus or the transform's correctness,
drives most of the difference.

The two reference sequences differ in two ways at once, and the
measurements above cannot separate them. The first is base composition.
The real 61-nucleotide sequence (Section~\ref{sec:refseq}) is close to
base-balanced: 15 A, 16 T, 15 C, 15 G, each base between 24.59\% and
26.23\% of the sequence. The 64-nucleotide sample sequence used in
Section~\ref{sec:revision-history} is visibly skewed: 20 A, 18 T, 16 C,
10 G, with A appearing exactly twice as often as G. Since A maps to the
value 0 (Section~\ref{sec:bytes-to-dna}), a reference that is
disproportionately A contributes nothing at a large fraction of its
positions, while a balanced reference contributes a varied offset
throughout.

The second is the period of the accumulated mask. By
Section~\ref{sec:mod4chain}, what the transform actually adds at
position $i$ is the running total $R_i$, whose period depends on
$S$, the sum of the reference sequence over one full cycle. The real
sequence gives $S = 91 \equiv 3 \pmod 4$, so $R$ has period
$4L = 244$. The 64-nucleotide sample sequence gives
$S = 80 \equiv 0 \pmod 4$, so its accumulated mask repeats every
$L = 64$ positions, exactly the periodicity the feedback term was
introduced to remove. The two candidate explanations therefore point
the same way for these two particular sequences, and no experiment
reported here distinguishes them. Section~\ref{sec:future-work}
proposes the sequences needed to separate the two effects.

For completeness: the real sequence contains the substring
\texttt{AGCT} five times (Section~\ref{sec:refseq}), which under the
mapping of Section~\ref{sec:bytes-to-dna} is the value sequence
$0,3,2,1$, one full cycle through all four mod-4 residues. This is a
coincidental feature of an already-balanced sequence, not evidence that
a repeated substring improves mixing: a repeated pattern is normally a
source of predictability, not a strength.

Entropy is also not the most sensitive signal available. Chi-square
shows a much larger relative separation from random, 4.8 to 8.4 times,
than entropy's under-1\% relative gap. A plausible reason is that
entropy is a compressed, log-scale summary that saturates near its
maximum quickly, while chi-square stays sensitive to smaller deviations
from uniformity for longer.

\section{Avalanche Effect}
\label{sec:results-avalanche}
Section~\ref{sec:measurement-design} explained why avalanche must be measured on
Pipeline A's \texttt{post\_dna} stage, not on the full pipeline's
AES-GCM ciphertext. This section reports the measured numbers and then
derives what they should have been.

Fig.~\ref{fig:avalanche-comparison} compares the two measurements
directly, for the three small corpus files. The full-pipeline
measurement (25 trials, from \texttt{results/experiment\_results.csv})
lands at 50.03\% (text), 50.00\% (json), and 50.04\% (random),
regardless of data class, confirming this number reflects AES-GCM's
random nonce and carries no information about the DNA transform. The
DNA-transform-only measurement, run with 2,000 trials, lands at 31.23\%
(text, stdev 1.66\%), 31.78\% (json, stdev 1.80\%), and 30.38\%
(random, stdev 1.79\%).

\begin{figure}[H]
\centering
\begin{tikzpicture}
\begin{axis}[
  ybar, bar width=16pt,
  width=11cm, height=6.5cm,
  ymin=0, ymax=62,
  ylabel={avalanche effect (\%)},
  y label style={font=\small},
  symbolic x coords={text,json,random},
  xtick=data,
  tick label style={font=\small},
  legend style={at={(0.5,-0.22)}, anchor=north, legend columns=1, font=\small},
  nodes near coords,
  every node near coord/.append style={font=\scriptsize},
  axis line style={gray!60},
]
\addplot+[fill=seqblue!45,draw=black!60] coordinates {(text,31.23)(json,31.78)(random,30.38)};
\addplot+[fill=seqblue!90,draw=black!60] coordinates {(text,50.03)(json,50.00)(random,50.04)};
\legend{{DNA transform only (2{,}000 trials, no nonce)},{Full pipeline (25 trials, includes AES-GCM nonce)}}
\end{axis}
\end{tikzpicture}
\caption{Avalanche effect: the percentage of output bits that flip when
one input bit is flipped, measured on Pipeline A's \texttt{post\_dna}
stage (2,000 trials, deterministic) versus the full pipeline's AES-GCM
ciphertext (25 trials). The full-pipeline bars land at essentially
50\% for every data class, reflecting AES-GCM's random nonce rather
than the DNA transform's own behaviour. The DNA-transform-only bars all
land near 31\%, the value predicted analytically in this section.}
\label{fig:avalanche-comparison}
\end{figure}

\textbf{Where 31\% comes from.} This figure is not evidence of
diffusion in the sense the metric was designed to capture. It follows
directly from the closed form in Equation~\ref{eq:prefixsum}. Flipping
one input bit changes exactly one input nucleotide $p_k$ by some
$\delta \ne 0$, and by Equation~\ref{eq:prefixsum} every output
nucleotide from position $k$ onward is shifted by that same constant
$\delta$, while every output nucleotide before position $k$ is
untouched.

Two quantities follow. First, how many bits a constant shift flips
within one affected nucleotide: for $\delta = 2$ the pair
$\{c, c+2\}$ always differs in exactly 1 of its 2 bits, so 50\%; for
$\delta = 1$ or $\delta = 3$ it differs in 1 or 2 bits depending on
$c$, averaging 1.5 of 2 bits, so 75\%. A randomly chosen input bit is
the high bit of its nucleotide (giving $\delta = 2$) or the low bit
(giving $\delta = 1$ or 3) with equal probability, so the mean over
affected nucleotides is $(50 + 75)/2 = 62.5\%$. Second, how much of the
output is affected at all: for a uniformly chosen flip position, half
the stream lies after it on average. The product,
$0.5 \times 62.5\% = 31.25\%$, is what the metric must report. The
three measured values, 30.38\% to 31.78\%, sit on that prediction.

Two consequences matter for how this number should be read. The spread
between the three data classes is not a finding: the derivation above
never refers to the plaintext, so the metric is essentially
data-independent by construction, and the differences between 30.38\%,
31.78\% and 31.23\% carry no interpretation. More importantly, the
transform has \emph{no backward diffusion at all}. Because the shift
applies only from the flip position onward, avalanche measured on this
transform is a linear function of where the bit was flipped: near 62.5\%
for a flip at the very start of the stream, and near 0\% for a flip at
the very end. A strong block cipher reports close to 50\% regardless of
flip position. Reporting a single mean over uniformly chosen positions,
as the standard one-bit-flip method does, hides this entirely.

The chained modulo-4 transform therefore propagates changes, but it does
not diffuse them in the sense a cipher does. Chapter~\ref{ch:limitations}
returns to what this implies about the transform's standing as a
security mechanism.

\section{Size Overhead}
\label{sec:results-overhead}
Section~\ref{sec:measurement-design} explained why \texttt{size\_overhead()},
computed against the AES-GCM ciphertext body alone, always returns
0.0\%. \texttt{results/table\_throughput\_size\_overhead.csv} confirms
this directly: the \texttt{size\_overhead\_body\_pct} column reads
exactly 0.0 for all 30 rows, without a single exception.

Counting the nonce and tag alongside the ciphertext, as
\texttt{size\_overhead\_full\_aead\_pct} does, gives a size-dependent number. Table~\ref{tab:overhead} reports it for all six
corpus files.

\begin{table}[h]
\centering
\caption{Full-AEAD size overhead (nonce + ciphertext + tag versus plaintext)}
\label{tab:overhead}
\begin{tabular}{lrr}
\toprule
File & Size (bytes) & Overhead (\%) \\
\midrule
text\_small.txt & 10238 & 0.3126 \\
json\_small.json & 10313 & 0.3103 \\
random\_small.bin & 10240 & 0.3125 \\
text\_large.txt & 204799 & 0.0156 \\
json\_large.json & 205012 & 0.0156 \\
random\_large.bin & 204800 & 0.0156 \\
\bottomrule
\end{tabular}
\end{table}

The fixed 32-byte overhead (a 16-byte nonce plus a 16-byte tag) costs
proportionally far more on the small files, about 0.31\%, than on the
large files, about 0.016\%, a difference of roughly 20 times. This
pattern is the same across all three data classes, since the overhead
depends only on the fixed nonce and tag sizes, not on the content of
the file.

% ============================================================
% CHAPTER 8 --- DISCUSSION
% ============================================================
\chapter{Discussion}
\label{ch:discussion}

\section{Answering the Research Question}
Chapter~\ref{ch:aim} stated two hypotheses. Hypothesis 1: the DNA-encoding
transform produces a measurably different statistical signature on
low-entropy data than on high-entropy data. Hypothesis 2: that effect does not
survive classical encryption.

H1 is supported, though the size of the effect depends on which metric
is used. Entropy alone shows only a small gap:
0.05 to 0.10 bits/byte between text or JSON and random data at the
\texttt{post\_dna} stage (Section~\ref{sec:results-correction}).
Chi-square shows a much clearer gap: text and JSON measure 4.8 to 8.4
times random's chi-square value at the same stage, or 38.3 and 75.5
standardised deviations from uniformity against random's $-0.9$
(Section~\ref{sec:results-entropy}), a separation confirmed at 48.1 and
46.9 standard errors across ten independent samples per class
(Section~\ref{sec:results-repeated-sample}).

The avalanche measurement does not contribute evidence for H1.
Section~\ref{sec:results-avalanche} shows the 30 to 32\% figure is the
analytically predicted consequence of the transform being a prefix sum,
independent of the plaintext, so it neither separates the data classes
nor demonstrates diffusion.

H2 is also supported. Every metric measured
at \texttt{post\_classical} lands within a narrow band, regardless of
data class or whether Pipeline A or B produced it
(Section~\ref{sec:results-entropy}). AES-GCM erases the DNA transform's
own statistical signature. Whatever the transform does to the data on
its own, it does not survive classical encryption.

\section{What the Entropy-Gap Correction Means}
Section~\ref{sec:results-correction} corrected an earlier, larger
estimate of the entropy gap, first reported in
Section~\ref{sec:revision-history}, down by roughly an order of
magnitude, and the cause is worth stating plainly. The gap's size was
not a fixed property of the DNA transform by itself. It depended on
which specific reference sequence was mixed with the data, and the two
sequences compared differ both in base composition and in the period of
the accumulated mask they induce. A result of this kind is therefore
not portable to a different reference sequence without re-measuring it,
and it is not yet known which of the two properties carries the effect.
Chapter~\ref{ch:limitations} returns to this point.

\section{Why Chi-Square Matters More Than Entropy Alone}
Entropy is the most familiar of the metrics used at the
\texttt{post\_dna} stage, but Section~\ref{sec:results-correction}
showed it is also the least sensitive one available here. A signal
that entropy alone would call negligible, under a 1\% relative
difference, still shows up clearly in chi-square, at 4.8 to 8.4 times
random's value and tens of standard deviations from uniformity. A
reader relying on entropy alone could reasonably conclude the DNA
transform has almost no effect on the data. That conclusion would be
wrong.

The avalanche metric illustrates the same lesson from the opposite
direction. Taken at face value, 30 to 32\% looks like evidence of
meaningful diffusion; derived from the transform's closed form, it
turns out to be a fixed consequence of the construction that says
nothing about the data at all (Section~\ref{sec:results-avalanche}).
Reporting several metrics side by side is what surfaces both cases, but
only if each metric's expected value under the null construction is
worked out rather than assumed.

\section{Practical Implications}
The size-overhead results
(Section~\ref{sec:results-overhead}) show AES-GCM's own on-wire cost is
small and shrinks further as messages grow, about 0.31\% for a 10 KB
message and 0.016\% for a 200 KB one. This cost comes from the
classical layer, not from the DNA transform: \texttt{post\_dna} output
is always the same length as the input
(Section~\ref{sec:bytes-to-dna}). The DNA transform's own computational
cost is a separate question, left open for the reasons
Section~\ref{sec:limits-performance} gives.

\section{Relation to DNA-LWCS and the Wider Literature}
Chapter~\ref{ch:related} noted that DNA-LWCS~\cite{aqeel2025} reports a
different set of metrics, correlation, encryption time, and energy
use, on image data, evaluated as a whole system against other published
schemes, so no direct numeric comparison with this thesis is possible.
What this thesis contributes instead is a question that complements
DNA-LWCS's own evaluation: whether a DNA-domain encoding step, on its
own, leaves a statistically detectable trace that depends on the
input's structure, and whether that trace survives once a real
classical cipher is layered on top. The answer here, yes to the first
and no to the second, is a kind of result an image-focused,
whole-system evaluation was not designed to produce.

% ============================================================
% CHAPTER 9 --- LIMITATIONS
% ============================================================
\chapter{Limitations}
\label{ch:limitations}

\section{Statistical Testing Is Not Cryptanalysis}
\label{sec:limits-not-cryptanalysis}
Chapter~\ref{ch:aim} already stated this limitation as a scope
boundary, adapted from NIST SP 800-22's own caveat~\cite{nist80022}:
passing or failing a statistical test says nothing directly about
whether a cipher, or a DNA-domain transform, is cryptographically
secure. This thesis's chi-square and entropy results
(Section~\ref{sec:results-entropy}) describe how uniform the byte
distribution looks. They do not demonstrate, or rule out, any specific
attack.

The transform's own structure makes the point concretely. By
Equation~\ref{eq:prefixsum}, consecutive output nucleotides satisfy
$c_i - c_{i-1} \equiv p_i + r_{i \bmod L} \pmod 4$, and the reference
sequence $r$ is public by design (Section~\ref{sec:refseq}). Anyone
holding \texttt{post\_dna} output can therefore recover the plaintext
exactly, in a single pass, by differencing adjacent nucleotides and
subtracting the known reference value. The flattened byte distribution
visible in Fig.~\ref{fig:byte-histogram} is cosmetic in this sense: the
original structure is fully present, one differencing step away. This
is not a flaw in the measurements, and it does not affect H1 or H2,
both of which concern statistical signatures rather than
recoverability. It does mean the \texttt{post\_dna} measurements should
be read as diagnostics of the transform and nothing more.
Confidentiality in the full pipeline rests on AES-GCM alone, and this
thesis makes no claim otherwise (Chapter~\ref{ch:intro}'s caveat).

\section{A Single, Fixed Reference Sequence}
Section~\ref{sec:results-correction} showed the measured entropy gap's
size depends heavily on which reference sequence is used, not only on
the transform's design. This thesis measured exactly one reference
sequence (Section~\ref{sec:refseq}): a 61-nucleotide sequence,
close to base-balanced, with two internal repeated runs, and with
$S \equiv 3 \pmod 4$. A different reference
sequence, shorter, longer, more skewed in its base composition, or with
a different value of $S$, could plausibly produce a smaller or larger
gap. This thesis does not test that range, and as
Section~\ref{sec:results-correction} notes, the one alternative
sequence it did compare against differs on both axes at once, so the
two candidate mechanisms remain confounded. The reference-sequence
dependency is itself a finding, but it also means the
exact gap sizes reported in Chapter~\ref{ch:results} should not be read
as a universal property of the chained modulo-4 design, only as this
specific sequence's result.

\section{Corpus Size and Composition}
The corpus covers three data classes at two sizes each, six files in
total (Section~\ref{sec:corpus}). This is enough to show a consistent
pattern across text, JSON, and random data, at a small and a large
size. It is not a large, statistically representative sample within
each class. The JSON files are synthetic log records, not captured
production traffic. The text files come from one 19th-century novel,
not a range of authors or writing styles. A wider corpus, with more
files per class, more classes, or real IoT traffic captures, could show
a different range of gap sizes than the ones reported here. The
underlying pattern itself, a signal at \texttt{post\_dna} that is
gone at \texttt{post\_classical}, held consistently across every file
actually tested, but six files is still a narrow base for a broader
claim.

\section{Implementation-Specific Performance}
\label{sec:limits-performance}
The DNA transform in this thesis is implemented as a direct,
pure-Python, per-nucleotide loop (Section~\ref{sec:module-structure}),
not an optimized or compiled implementation. Any conclusion about the
transform's computational cost would depend heavily on implementation
choices not tested here.
This thesis's results chapter reports statistical signal, not detailed
timing comparisons, precisely to avoid overstating what one specific,
unoptimized implementation's speed says about the DNA-domain approach
in general. Chapter~\ref{ch:intro}'s caveat, that this thesis does not
test whether DNA encoding is lightweight in a performance sense,
follows directly from this.

\section{A Non-Standard AES-GCM Nonce Length}
\label{sec:limits-nonce}
The size-overhead figures in Section~\ref{sec:results-overhead} assume
PyCryptodome's default AES-GCM nonce length of 16 bytes (128 bits),
which this project uses rather than explicitly requesting a shorter
one. Both standards cited in this thesis point elsewhere. NIST SP
800-38D recommends a 96-bit (12-byte) nonce and specifies a different,
GHASH-based derivation of the internal counter block for any other
length~\cite{nist80038d}. RFC 5116 goes further for the registered
algorithm: AEAD\_AES\_256\_GCM fixes both the minimum and the maximum
nonce length at 12 octets~\cite{rfc5116}, so a 16-byte nonce places
this project's construction outside the algorithm whose 16-octet
ciphertext expansion Section~\ref{sec:measurement-design} cites.

The consequence for the reported numbers is small: the fixed overhead
is 32 bytes (16-byte nonce plus 16-byte tag) rather than the 28 bytes
a 12-byte nonce would give, moving a 10 KB file's overhead from
0.3126\% to about 0.2735\%. That is far smaller than the roughly
20-fold gap between small- and large-file overhead on which
Section~\ref{sec:results-overhead}'s conclusion actually rests. It is
nonetheless an unrequested departure from both cited standards rather
than a tested design choice, and a corrected implementation should
request a 12-byte nonce explicitly.

\section{Scope Boundaries Not Revisited}
Two exclusions announced in Chapter~\ref{ch:aim} are worth restating as
limitations rather than choices. ECC/ECDH was not implemented, so the
original DNA-plus-ECC combination was never run even once as a sanity
check against this thesis's AES-GCM baseline; the ablation in
Fig.~\ref{fig:pipelines} does not depend on it, but neither does this
thesis verify that its DNA-domain stage behaves as the source works'
does. The full NIST STS suite was likewise not run, so the statistical
picture rests on the five metrics of Section~\ref{sec:metrics} rather
than on a standard battery.

% ============================================================
% CHAPTER 10 --- SUMMARY AND FUTURE WORK
% ============================================================
\chapter{Summary and Future Work}
\label{ch:summary}

\section{Contributions}
This thesis set out to determine whether a DNA-domain encoding step
leaves a measurable, structure-dependent statistical trace on data,
and whether that trace survives AES-GCM encryption. Concretely, this
work produced the following:

\begin{itemize}
\item \textbf{An invertible DNA-domain transform with feedback.} The
chained modulo-4 design (Section~\ref{sec:mod4chain}) combines a
public reference sequence with the plaintext and a feedback term. Its
round-trip test suite confirms it correctly recovers the original data
across edge cases, including random-length inputs, empty input, every
possible byte value, and text and binary samples
(\texttt{tests/test\_dna\_transform.py}).

\item \textbf{A closed-form characterisation of that transform.}
Equation~\ref{eq:prefixsum} shows the design is a prefix sum, which
determines both its avalanche behaviour
(Section~\ref{sec:results-avalanche}) and the period of the mask it
applies (Section~\ref{sec:refseq}).

\item \textbf{An ablation methodology isolating the transform's
effect.} Two pipelines share the same key and cipher and differ only
in whether the DNA step runs (Section~\ref{sec:pipeline-design}).

\item \textbf{A reproducible, checksummed test corpus.} Six frozen
files across three data classes and two sizes
(Section~\ref{sec:corpus}), with entropy, chi-square, compression,
correlation, avalanche, and overhead measured at each pipeline stage.

\item \textbf{A modular, tested implementation.} The pipeline is split
into single-purpose modules, each with its own test file
(Section~\ref{sec:module-structure}).

\item \textbf{A measurable pre-encryption signal.} Text and JSON
produce a chi-square statistic 4.8--8.4 times higher than random
input at the \texttt{post\_dna} stage, separated by 48.1 and 46.9
standard errors across ten samples per class
(Sections~\ref{sec:results-entropy}, \ref{sec:results-repeated-sample}).

\item \textbf{No surviving signal after AES-GCM.} Every metric
measured after encryption lands within a narrow band, regardless of
data class or whether the DNA step ran
(Section~\ref{sec:results-entropy}).

\item \textbf{A documented account of two rejected designs.} A
SHA-256-mask-and-XOR transform erased the entropy difference the
experiment depends on measuring
(Section~\ref{sec:revision-history}). A tiled mask without feedback
preserved that difference but produced an autocorrelation of 0.835 at
the reference sequence's own length, and was rejected for it.

\end{itemize}

A DNA-domain encoding step of this kind is not a substitute for
classical encryption. It has a small statistical effect of its own
before encryption, and AES-GCM removes that effect completely.

\section{Future Work}
\label{sec:future-work}
Four items follow directly from the limitations above, in rough order
of how cheaply they could be done.

\textbf{Separate base balance from mask period.}
Section~\ref{sec:results-correction} leaves two candidate explanations
for the reference sequence's effect confounded, because the two
sequences compared differ on both. Constructing four reference
sequences that cross the two factors, base-balanced and skewed against
$S \equiv 0$ and $S \not\equiv 0 \pmod 4$, and re-measuring the
entropy gap for each would resolve which one carries the effect. The
same sweep would test whether the lag-$L$ autocorrelation that
Design 2 suffered from returns when $S \equiv 0$, as
Section~\ref{sec:mod4chain} predicts it should.

\textbf{Re-measure the swap analysis under a committed script.} The
reference-sequence swaps in Section~\ref{sec:results-correction} were
partly ad hoc, and one of them is not reproducible. Persisting all of
them as a script under \texttt{scripts/} would bring that analysis up
to the reproducibility standard the rest of the project follows, and
would also cover the Design 3 autocorrelation re-measurement noted in
Section~\ref{sec:revision-history}.

\textbf{Report avalanche as a function of flip position.}
Section~\ref{sec:results-avalanche} shows the single averaged avalanche
figure conceals the transform's lack of backward diffusion. Plotting
Hamming distance against the position of the flipped bit would make
that visible in one figure, and would provide a template for evaluating
any other stream-structured transform against the same metric.

\textbf{Extend the statistical battery and the corpus.} The five
metrics used here could be extended toward the NIST STS suite, in
particular the serial and runs tests, which are sensitive to the kind
of positional structure a prefix sum leaves behind and which
Pearson correlation demonstrably is not
(Section~\ref{sec:results-correlation}). Pairing that with captured IoT
traffic rather than synthetic JSON would address the corpus
limitation directly.

% ============================================================
% BIBLIOGRAPHY
% ============================================================
\begin{thebibliography}{99}

\bibitem{iotanalytics2025}
IoT Analytics,
\emph{Number of Connected IoT Devices Growing 14\% to 21.1 Billion
Globally}, 2025.
\url{https://iot-analytics.com/number-connected-iot-devices/}

\bibitem{sonicwall2024}
SonicWall,
\emph{2024 SonicWall Mid-Year Cyber Threat Report}, 2024.

\bibitem{stationx2026}
StationX,
\emph{IoT Security Statistics: Threats and Trends}, 2026.
\url{https://www.stationx.net/iot-security-statistics/}

\bibitem{unit42-2020}
Palo Alto Networks, Unit 42,
\emph{2020 Unit 42 IoT Threat Report}, 2020.
\url{https://unit42.paloaltonetworks.com/iot-threat-report-2020/}

\bibitem{zinabu2025}
Zinabu N.~G., Marye Y.~W., Tune K.~K., Demilew S.~A.,
\emph{Comprehensive Analysis of Lightweight Cryptographic Algorithms
for Battery-Limited Internet of Things Devices}, International Journal
of Distributed Sensor Networks, 2025.
\url{https://doi.org/10.1155/dsn/9639728}

\bibitem{thakor2021}
Thakor V.~A., Razzaque M.~A., Khandaker M.~R.~A.,
\emph{Lightweight Cryptography Algorithms for Resource-Constrained IoT
Devices: A Review, Comparison and Research Opportunities}, IEEE Access,
vol.~9, pp.~28177--28193, 2021.
\url{https://doi.org/10.1109/ACCESS.2021.3052867}

\bibitem{nist800232}
Turan M.~S., McKay K., et~al.,
\emph{Ascon-Based Lightweight Cryptography Standards for Constrained
Devices}, NIST Special Publication 800-232, National Institute of
Standards and Technology, 2025.
\url{https://doi.org/10.6028/NIST.SP.800-232}

\bibitem{aqeel2025}
Aqeel S., Khan A.~S., Abbasi I.~A., Algarni F., Grzonka D.,
\emph{Enhancing IoT security with a DNA-based lightweight cryptography
system}, Scientific Reports, vol.~15, art.~13367, 2025.
\url{https://doi.org/10.1038/s41598-025-96292-0}

\bibitem{dnaecc}
Barman P., Saha B.,
\emph{DNA Encoded Elliptic Curve Cryptography System for IoT Security},
arXiv:2311.11393, 2023.
\url{https://arxiv.org/abs/2311.11393}

\bibitem{nist80022}
Rukhin A., et~al.,
\emph{A Statistical Test Suite for Random and Pseudorandom Number
Generators for Cryptographic Applications}, NIST Special Publication
800-22 Rev.~1a, National Institute of Standards and Technology, 2010.
\url{https://doi.org/10.6028/NIST.SP.800-22r1a}

\bibitem{nist80038d}
Dworkin M.,
\emph{Recommendation for Block Cipher Modes of Operation:
Galois/Counter Mode (GCM) and GMAC}, NIST Special Publication
800-38D, National Institute of Standards and Technology, 2007.
\url{https://doi.org/10.6028/NIST.SP.800-38D}

\bibitem{sec1}
Certicom Research,
\emph{SEC 1: Elliptic Curve Cryptography}, Standards for Efficient
Cryptography, Version 2.0, 2009.
\url{https://www.secg.org/sec1-v2.pdf}

\bibitem{inru}
Tiwari S.~K., Awasthi A., Chakrabarti S., Yadav S.,
\emph{INRU: A Quasigroup Based Lightweight Block Cipher},
arXiv:2112.07411, 2021.
\url{https://arxiv.org/abs/2112.07411}

\bibitem{ascon-survey}
Kaur J., Cintas-Canto A., Mozaffari-Kermani M., Azarderakhsh R.,
\emph{A Comprehensive Survey on the Implementations, Attacks, and
Countermeasures of the Current NIST Lightweight Cryptography Standard
(ASCON)}, arXiv:2304.06222, 2023.
\url{https://arxiv.org/abs/2304.06222}

\bibitem{rfc5116}
McGrew D.,
\emph{An Interface and Algorithms for Authenticated Encryption}, RFC
5116, Internet Engineering Task Force, 2008.
\url{https://www.rfc-editor.org/rfc/rfc5116.html}

\end{thebibliography}

% TODO: the bibliography is still thin for a Master's thesis (15 entries).
% Related Work rests on two DNA-cryptography papers, one of which
% Section 4.1 admits was not reviewed in depth. Broaden with:
%   - the wider DNA-cryptography literature (Gehani/Reif/Bancroft and
%     the DNA image-encryption line that followed)
%   - additional lightweight-cipher benchmarking work beyond [ascon-survey]
%   - a source for the classical entropy of written English
% See docs/milestone_plan.md.

% ============================================================
% LIST OF FIGURES / TABLES / LISTINGS
% ============================================================
\listoffigures
\addcontentsline{toc}{chapter}{List of Figures}
\listoftables
\addcontentsline{toc}{chapter}{List of Tables}
\lstlistoflistings
\addcontentsline{toc}{chapter}{List of Listings}

\end{document}